\documentclass[a4paper,12pt]{article}
%\usepackage[lite,subscriptcorrection,slantedGreek,nofontinfo]{mtpro2}
%\usepackage{times} 
\usepackage{amsfonts}
\usepackage{mathrsfs}
\usepackage{amsmath}
\usepackage{amssymb}
\usepackage{framed}
%\usepackage{mathptmx}
\usepackage[medium]{titlesec}
\usepackage{bm}
\usepackage{cite}                                                   
%\usepackage{feynmp-auto}
\usepackage[normalem]{ulem}
\usepackage{extarrows}
\usepackage{slashed}
\usepackage{isodateo}
\usepackage{graphicx}
\usepackage{xcolor}
\usepackage[bookmarksnumbered=true,bookmarksopen=true]{hyperref}
 \hypersetup{colorlinks,%
             linkcolor=[rgb]{0,0.3,0.6}, %
             citecolor=[rgb]{0,0.3,0.6}, %
             urlcolor=[rgb]{0,0.3,0.6}}
\usepackage[hmargin=.7in,vmargin=1.1in]{geometry}
\usepackage{indentfirst}
\usepackage{booktabs}

\linespread{1.1}
%\numberwithin{equation}{section}
\newcommand{\FR}[2]{\displaystyle\frac{\,{#1}\,}{#2}}
\newcommand{\fr}[2]{\mbox{$\frac{\,{#1}\,}{#2}$}}
\newcommand{\n}{\nonumber}
\renewcommand{\rm}{\mathrm}
\renewcommand{\thefootnote}{\arabic{footnote}}

\newcommand{\red}{\color{red}}
\newcommand{\blue}{\color{blue}}
\newcommand{\gray}{\color{gray}}
\newcommand{\hl}{\color[rgb]{1,0,.5}}

\def\bge{\begin{equation}}
\def\ede{\end{equation}}
\def\bga{\begin{aligned}}
\def\eda{\end{aligned}}
\def\bgb{\begin{bmatrix}}
\def\edb{\end{bmatrix}}
\def\bgp{\begin{pmatrix}}
\def\edp{\end{pmatrix}}
\def\bgm{\begin{matrix}}
\def\edm{\end{matrix}}
\def\bgs{\begin{subequations}}
\def\eds{\end{subequations}}
\newcommand{\order}[1]{\mathcal{O}({#1})}
\def\di{{\mathrm{d}}}
\def\D{{\mathrm{D}}}
\def\Di{{\mathcal{D}}}

\def\T{\mathcal{T}}
\def\J{\mathcal{J}}
\def\mb{\mathbf}
\def\ms{\mathscr}
\def\mf{\mathfrak}
\def\mc{\mathcal}
\def\cl{\mathrm{cl}}
\def\pd{\partial}
\def\ld{{\mathscr{L}}}
\def\hd{{\mathscr{H}}}
\def\z{{\bar{z}}}
\def\la{\langle}\def\ra{\rangle}
\def\sla{\slashed}
\def\const{\mathrm{const.}}
\setlength\unitlength{1mm}
\def\tr{\mathrm{\,tr\,}}
\def\Tr{\mathrm{\,Tr\,}}
\def\Det{\mathrm{\,Det\,}}
\def\to{\rightarrow}
\def\To{\Rightarrow}
\def\ii{\mathrm{i}}

\def\al{\alpha}
\def\be{\beta}
\def\ga{\gamma}
\def\de{\delta}
\def\ep{\epsilon}
\def\ka{\kappa}
\def\lam{\lambda}
\def\rh{\rho}
\def\si{\sigma}
\def\ze{\zeta}
\def\tn{\tilde\nu}


\def\Wh{\mathrm{W}}


\def\aa{\mathsf{a}}
\def\bb{\mathsf{b}}
\def\cc{\mathsf{c}}
\def\dd{\mathsf{d}}

\def\C{\mathbb{C}}
\def\R{\mathbb{R}}

\def\Mp{M_{\text{Pl}}}


\def\Re{\mathrm{Re}\,}
\def\Im{\mathrm{Im}\,}

\def\ad{\mathrm{\,ad\,}}

\usepackage{mdframed} 
                   

\newmdenv[skipabove=0pt,%
          skipbelow=5pt,%
          leftmargin=0pt,%
          rightmargin=0pt,%
          innertopmargin=-5pt,%
          innerbottommargin=7pt,%
          innerleftmargin=2pt,%
          innerrightmargin=2pt,%
          splittopskip=0pt,%
          splitbottomskip=0pt,%
          linewidth=0pt,%
          nobreak=true]%
          {keyeqn2}                   
                   
                   
\newmdenv[backgroundcolor=gray!15,%
          skipabove=0pt,%
          skipbelow=5pt,%
          leftmargin=0pt,%
          rightmargin=0pt,%
          innertopmargin=-5pt,%
          innerbottommargin=7pt,%
          innerleftmargin=2pt,%
          innerrightmargin=2pt,%
          splittopskip=0pt,%
          splitbottomskip=0pt,%
          linewidth=0pt,%
          nobreak=true]%
          {keyeqn}

\usepackage{titlesec}          
\titleformat{\section}
{\normalfont\fontsize{15}{20}\bfseries}{\thesection}{1em}{}

\usepackage[T1]{fontenc}
\usepackage{lmodern}

\newcommand{\ob}[1]{\mkern 2mu \overline{\mkern -2mu #1 \mkern -2mu}\mkern 2mu}
\newcommand{\wt}[1]{\mkern 2mu \widetilde{\mkern -2mu #1 \mkern -2mu}\mkern 2mu}
\newcommand{\wh}[1]{\mkern 2mu \widehat{\mkern-2mu#1\mkern-2mu}\mkern 2mu}

\newcommand{\fnemail}[1]{\footnote{Email: \href{mailto:#1}{\nolinkurl{#1}}}}

\title{\Large\textbf{Standard Model Signal in Cosmological Higgs Collider}\\[2mm]}

\author{}

\begin{document}

\date{}
\maketitle

% \vspace{20mm}

% \begin{abstract}
% \vspace{10mm}

% We derive ...

% \end{abstract}

% % \newpage
% \tableofcontents

% \newpage
% 3 methods: Late-time expansion of two point function; SK integral in time-momentum space and factorization theorem; spectral decomposition. The bubble diagram can be evaluated in these 3 methods, while the triangle diagrams and box diagram can only be evaluated by the factorization theorem method.
\section{Gauge Boson Signal}
Higgs bulk-boundary propagator is a massless propagator:
\begin{align}
    G_\pm(k; \tau) = \frac{H^2}{2k^3} (1 \mp ik\tau) e^{\pm ik\tau}
\end{align}

The Higgs-gauged coupling are from the original SM Lagrangian: 
\begin{align}
    \frac{\sqrt{-g}}{2}|\mathrm{D}_\mu \mb H |^2 \supset \frac{ a(\tau)^2}{4} g^2 (h_0 + \de h)^2 \left( W^+_\mu W^{-\mu} + \frac{1}{2c_W^2} Z^\mu Z_\mu \right)
\end{align}

The masses of gauge bosons due to the infrared Higgs loop contribution are:
\begin{align}
    m_W^2 = \frac{3 g^2 H^4}{8 \pi^2 m_H^2}, \quad m_Z^2 = \frac{m^2_W}{c^2_W}
\end{align}

Mode function of a massive gauge boson \cite{qin_helical_2023}:
\begin{align}
    B^{( \pm)}(k, \tau) & =\frac{1}{\sqrt{2 k}} \mathrm{~W}_{ 0, \mathrm{i} \tilde{\nu}}(2 \mathrm{i} k \tau), \\
    B^{(T)}(k, \tau) & =\frac{\sqrt{\pi} k}{2 m} e^{-\pi \tilde{\nu} / 2} H(-\tau)^{3 / 2} \mathrm{H}_{\mathrm{i} \tilde{\nu}}^{(1)}(-k \tau), \\
    B^{(L)}(k, \tau) & =-\frac{\mathrm{i} \sqrt{\pi}}{4 m} e^{-\pi \widetilde{\nu} / 2} H(-\tau)^{1 / 2}\left[\mathrm{H}_{\mathrm{i} \tilde{\nu}}^{(1)}(-k \tau)-k \tau\left(\mathrm{H}_{1+\mathrm{i} \tilde{\nu}}^{(1)}(-k \tau)-\mathrm{H}_{-1+\mathrm{i} \tilde{\nu}}^{(1)}(-k \tau)\right)\right]
\end{align}
where $\tilde\nu=\sqrt{(m/H)^2-1/4}$. And the W function with zero chemical potential can transform to a Hankel function:
\begin{align}
    \mathrm{~W}_{ 0, \mathrm{i} \tilde{\nu}}(2 \mathrm{i} k \tau) &= \sqrt{\frac{2ik\tau}{\pi}} K_{i \tilde\nu}(ik\tau) \nonumber \\
    &= \sqrt{\frac{ i\pi k\tau}{2}} ie^{-\pi \tilde\nu /2} \mathrm{H}^{(1)}_{i\tilde\nu}(-k\tau)
\end{align}

So the grater propagator or $+-$ propagator of each polarization:
\begin{align}
    &D_{>}^{( \pm)}\left(k ; \tau_{1}, \tau_{2}\right)=B^{( \pm)}\left(k, \tau_{1}\right) B^{( \pm) *}\left(k, \tau_{2}\right)=\frac{\pi e^{-\pi\tilde\nu} (\tau_1 \tau_2)^{1/2}}{4} \mathrm{H}_{\mathrm{i} \widetilde{\nu}}^{(1)}\left(-k \tau_{1}\right) \mathrm{H}_{-\mathrm{i} \widetilde{\nu}}^{(2)}\left(-k \tau_{2}\right), \\
    &D_{>}^{(T)}\left(k ; \tau_{1}, \tau_{2}\right)=B^{(T)}\left(k, \tau_{1}\right) B^{(T) *}\left(k, \tau_{2}\right)=\frac{\pi H^{2} k^{2}}{4 m^{2}} e^{-\pi \widetilde{\nu}}\left(\tau_{1} \tau_{2}\right)^{3 / 2} \mathrm{H}_{\mathrm{i} \widetilde{\nu}}^{(1)}\left(-k \tau_{1}\right) \mathrm{H}_{-\mathrm{i} \widetilde{\nu}}^{(2)}\left(-k \tau_{2}\right), \\
    &D_{>}^{(L)}\left(k ; \tau_{1}, \tau_{2}\right)=\frac{1}{k^{2}}\left(\partial_{\tau_{1}}-\frac{2}{\tau_{1}}\right)\left(\partial_{\tau_{2}}-\frac{2}{\tau_{2}}\right) D_{>}^{(T)}\left(k ; \tau_{1}, \tau_{2}\right) .
\end{align}

We quote the scalar propagator with mass index $i\tilde\nu$:
\begin{align}
    D^{(\text{scalar})}_> = \frac{\pi }{4 } e^{-\pi \widetilde{\nu}}\left(\tau_{1} \tau_{2}\right)^{3 / 2} \mathrm{H}_{\mathrm{i} \widetilde{\nu}}^{(1)}\left(-k \tau_{1}\right) \mathrm{H}_{-\mathrm{i} \widetilde{\nu}}^{(2)}\left(-k \tau_{2}\right)
\end{align}
and the nonlocal part of its late time expansion:
\begin{align}
    [D_>]_{\text{NonL}} = [D_<]_{\text{NonL}} = \frac{(\tau_1 \tau_2)^{3/2}}{4\pi} \left[ \Gamma^2(-i\tn) \left( \frac{k^2\tau_1\tau_2}{4} \right)^{i\tn} +\text{c.c.} \right]
\end{align}

And the vector propagator can be related to the scalar propagator:
\begin{align}
    D_>^{(\pm)} &= \frac{1}{\tau_1 \tau_2} D^{(\text{scalar})}_> \\
    D_>^{(T)} &= \frac{k^2}{\tilde\nu^2 + 1/4} D^{(\text{scalar})}_> \\
    D_{>}^{(L)} &= \frac{1}{\tilde\nu^2 + 1/4}\left(\partial_{\tau_{1}}-\frac{2}{\tau_{1}}\right)\left(\partial_{\tau_{2}}-\frac{2}{\tau_{2}}\right) D_{>}^{(\text{scalar})}\left(k ; \tau_{1}, \tau_{2}\right) 
\end{align}
We can see the propagator of temporal component of longitudinal mode have a suppression factor $(k\tau)^2$, so it will contribute the next leading term and we will ignore it below.

And the nonlocal part:
\begin{align}
    [D^{(\pm)}]_{\text{NonL}} &= \frac{(\tau_1 \tau_2)^{1/2}}{4\pi} \left[ \Gamma^2(-i\tn) \left( \frac{k^2\tau_1\tau_2}{4} \right)^{i\tn} +\text{c.c.} \right] \\
    [D^{(T)}]_{\text{NonL}} &= \frac{k^2}{\tilde\nu^2 + 1/4} \frac{(\tau_1 \tau_2)^{3/2}}{4\pi} \left[ \Gamma^2(-i\tn) \left( \frac{k^2\tau_1\tau_2}{4} \right)^{i\tn} +\text{c.c.} \right] \\
    [D^{(L)}]_{\text{NonL}} &= \frac{(-1/2+i\tn)^2}{\tilde\nu^2 + 1/4} \frac{(\tau_1 \tau_2)^{1/2}}{4\pi} \left[ \Gamma^2(-i\tn) \left( \frac{k^2\tau_1\tau_2}{4} \right)^{i\tn} +\text{c.c.} \right]
\end{align}
% We can see that the temporal component have a suppression of factor $(\tau_1\tau_2)$, so it will contribute the next leading term and we will neglected it in the following calculation.

We also define the propagator of helicity eigenstates:
\begin{align}
    &D_{>\;\mu\nu}^{{(\pm1)}} \left(\mb k ; \tau_{1}, \tau_{2}\right) = e_{\mu,\mb k}^{(\pm) } e_{\nu, \mb k}^{(\pm)*} D_{>}^{\pm} \left(k ; \tau_{1}, \tau_{2}\right) \\
    &D_{>\;\mu\nu}^{{(0)}} \left(\mb k ; \tau_{1}, \tau_{2}\right) =  e_{\mu,\mb k}^{(L)} e_{\nu, \mb k}^{(L)*} D_{>}^{(L)}\left(k ; \tau_{1}, \tau_{2}\right) \\
    &D_{>\;\mu\nu}^{} \left(\mb k ; \tau_{1}, \tau_{2}\right) = \sum_{h=\pm1,0} D_{>\;\mu\nu}^{{(h)}} \left(k ; \tau_{1}, \tau_{2}\right)
\end{align}

For $\mb{k}  = (0,0,k)$, the polarization vector is:
\begin{align}
    e_{\mu,\mathbf{k}}^{(T)} = \begin{pmatrix} 1 \\ 0 \\ 0 \\ 0 \end{pmatrix} , \qquad
    e_{\mu,\mathbf{k}}^{(\pm)} = \frac{1}{\sqrt{2}} \begin{pmatrix} 0 \\ 1 \\ \pm \mathrm{i} \\ 0 \end{pmatrix} , \qquad
    e_{\mu,\mathbf{k}}^{(L)} = \begin{pmatrix} 0 \\ 0 \\ 0 \\ 1 \end{pmatrix}, \qquad
    k_\mu = \begin{pmatrix} 0 \\ 0 \\ 0 \\ k \end{pmatrix}, \qquad
\end{align}
and when we rotate $\mb{{k}}$, the $e_{\mu,\mathbf{k}}^{(T)}$ is invariant and the other two types of porlarization vector will transform $e_{i,\mathbf{k}}' = R_{ij} e_{j,\mathbf{k}} $:
\begin{align}
    e_{\mu,\mathbf{k}}^{(\pm)} = \frac{1}{\sqrt{2}} \begin{pmatrix} 0 \\ \cos\theta\cos\phi \mp \mathrm{i}\sin\phi \\ \cos\theta\sin\phi \pm \mathrm{i}\cos\phi \\ -\sin\theta \end{pmatrix}, \qquad
    k_\mu = e_{\mu,\mathbf{k}}^{(L)} = \begin{pmatrix} 0 \\ \sin\theta\cos\phi \\ \sin\theta\sin\phi \\ \cos\theta \end{pmatrix}
\end{align}


The MB representation of the scalar propagator:
\begin{align}
    D_\gtrless \left(k ; \tau_{1}, \tau_{2}\right)= & \frac{1}{4 \pi} \int_{-\mathrm{i} \infty}^{+\mathrm{i} \infty} \frac{\mathrm{~d} s_{1}}{2 \pi \mathrm{i}} \frac{\mathrm{~d} s_{2}}{2 \pi \mathrm{i}} e^{\pm \mathrm{i} \pi\left(s_{1}-s_{2}\right)}\left(\frac{k}{2}\right)^{-2 s_{12}}\left(-\tau_{1}\right)^{-2 s_{1}+3 / 2}\left(-\tau_{2}\right)^{-2 s_{2}+3 / 2} \\
    & \times \Gamma\left[s_{1}-\frac{\mathrm{i} \widetilde{\nu}}{2}, s_{1}+\frac{\mathrm{i} \widetilde{\nu}}{2}, s_{2}-\frac{\mathrm{i} \widetilde{\nu}}{2}, s_{2}+\frac{\mathrm{i} \widetilde{\nu}}{2}\right]
\end{align}

For simplicity we define some loop seeds:
\begin{align}
    &\int \frac{d^3\mathbf{q}}{(2\pi)^3} |\mathbf{q}|^{-2a} |\mathbf{q} + \mathbf{k}_s|^{-2b} = \frac{k_s^{3-2a-2b}}{(4\pi)^{3/2}} \cdot f_0(a,b) \\
    &\int \frac{d^3\mathbf{q}}{(2\pi)^3} |\mathbf{q}|^{-2a} |\mathbf{q} + \mathbf{k}_s|^{-2b} \cos\theta = \frac{k_s^{3-2a-2b}}{(4\pi)^{3/2}} \cdot f_c(a,b) \\
    &\int \frac{d^3\mathbf{q}}{(2\pi)^3} |\mathbf{q}|^{-2a} |\mathbf{q} + \mathbf{k}_s|^{-2b} \cos^2\theta = \frac{k_s^{3-2a-2b}}{(4\pi)^{3/2}} \cdot f_{cc}(a,b) 
\end{align}
where $f(a,b)$ are some Gamma functions arising from the loop integral:
\begin{align}
    &f_0(a,b) = \Gamma \left[\begin{array}{c}
    a+b-\frac{3}{2}, \frac{3}{2}-a, \frac{3}{2}-b \\
    3-a-b, a,b
    \end{array}\right] \\
    &f_c(a,b) = \frac{1}{2}f_0(a+1/2,b-1) - \frac{1}{2}f_0(a-1/2,b) - \frac{1}{2}f_0(a+1/2,b) \\
    &f_{cc}(a,b) = \frac{1}{4}f_0(a+1,b-2) + \frac{1}{4}f_0(a-1,b) + \frac{1}{4}f_0(a+1,b) \nonumber \\
    &\qquad \quad - \frac{1}{2}f_0(a,b-1) -\frac{1}{2}f_0(a+1,b-1) +\frac{1}{2}f_0(a,b)
\end{align}

\section{Bubble Diagram}
For the Bubble diagram:
\begin{align}
    \la \delta h^4 \ra^{\text{Bubble}} &= \left( g^2 \right)^2 \sum_{a,b=\pm} (-ab) \int \frac{\di \tau_1}{|H \tau_1|^2} \frac{\di \tau_2}{|H \tau_2|^2} G_a(k_1;\tau_1) G_a(k_2;\tau_1) G_b(k_3;\tau_2) G_b(k_4;\tau_2) \nonumber \\
    &\times \int \frac{\di \mb{q}^3}{(2\pi)^3} \left[ D_{\mu\nu}(q,\tau_1,\tau_2) \right]_{ab} \left[ D^{\nu\mu}(|\mb{k}_s + \mb{q}|,\tau_2,\tau_1) \right]_{ba}
\end{align}

The loop integral part is:
\begin{align}
    \mathcal{LOOP}_{ab} &= \int \frac{\di \mb{q}^3}{(2\pi)^3} \left[ D_{\mu\nu}(q,\tau_1,\tau_2) \right]_{ab} \left[ D^{\nu\mu}(|\mb{k}_s + \mb{q}|,\tau_2,\tau_1) \right]_{ba}
\end{align}

Using the cutting rule, we know that the nonlocal signal part is independent of the SK index:
\begin{align}
    \mathcal{LOOP} = \int \frac{\di \mb{q}^3}{(2\pi)^3} &\left[ D^{(L)}(q,\tau_1,\tau_2) e^{(L)i}_{\mb q} e^{(L)j*}_{\mb q} + D^{(\pm)}(q,\tau_1,\tau_2) e^{(\pm)i}_{\mb q} e^{(\pm)j*}_{\mb q} \right] \nonumber \\
    \times &\left[ D^{(L)}(K,\tau_2,\tau_1) e^{(L)j}_{\mb K} e^{(L)i*}_{\mb K} + D^{(\pm)}(K,\tau_2,\tau_1) e^{(\pm)j}_{\mb K} e^{(\pm)i*}_{\mb K} \right] \\
    &= \sum_{\sigma_1,\sigma_2=L,\pm} \mathcal{{LOOP}}_{\sigma_1\sigma_2}
\end{align}
where $\mb K = \mb q+\mb k_s$. So we should deal with different polarization mixing. We set $\mb k_s=k_s(0,0,1)$ and $\mb q = q(\sin\theta\sin\phi,\sin\theta\cos\phi,\cos\theta)$, then the polarization vector become:
\begin{align}
    e_{i,\mathbf{q}}^{(L)} = \begin{pmatrix} \sin\theta\cos\phi \\ \sin\theta\sin\phi \\ \cos\theta \end{pmatrix},\qquad
    e_{i,\mathbf{q}}^{(\pm)} = \frac{1}{\sqrt{2}} \begin{pmatrix} \cos\theta\cos\phi \mp \mathrm{i}\sin\phi \\ \cos\theta\sin\phi \pm \mathrm{i}\cos\phi \\ -\sin\theta \end{pmatrix}
\end{align}

\begin{align}
    \mathbf{e}_{i,\mathbf{K}}^{(L)} = 
    \frac{1}{K}
    \begin{pmatrix}
    q\sin\theta\cos\phi \\
    q\sin\theta\sin\phi \\
    q\cos\theta + k_s
    \end{pmatrix},\qquad
    \mathbf{e}_{i,\mathbf{K}}^{(\pm)} = 
    \frac{1}{\sqrt{2} K}
    \begin{pmatrix}
    \cos\phi\,(q\cos\theta + k_s) \mp i\,\sin\phi\,K \\
    \sin\phi\,(q\cos\theta + k_s) \pm i\,\cos\phi\,K\\
    -q\sin\theta
    \end{pmatrix}
\end{align}

The signal part is contributed from the loop integral of nonlocal part of each propagator. And we calculate the first term as an example:
\begin{align}
    \mathcal{{LOOP}}_{LL}
    &=\int \frac{\di \mb{q}^3}{(2\pi)^3} \left[ D^{(L)}(q,\tau_1,\tau_2) e^{(L)i}_{\mb q} e^{(L)j*}_{\mb q} \cdot D^{(L)}(K,\tau_2,\tau_1) e^{(L)j}_{\mb K} e^{(L)i*}_{\mb K} \right]_{\text{NonL}} \nonumber \\
    &= \sum_{c=\pm} \frac{(-1/2+i\tn)^4}{(\tilde\nu^2 + 1/4)^2} \frac{(\tau_1 \tau_2)^{1+2ic\tn}}{(4\pi)^2} \Gamma^4(-ic\tn) \int \frac{\di \mb{q}^3}{(2\pi)^3}\left[  \left( \frac{q}{2} \right)^{2ic\tn} \left( \frac{K}{2} \right)^{2ic\tn} |e^{(L)}_{\mb q} \cdot e^{(L)*}_{\mb K}|^2 \right] 
\end{align}

Completing the loop integral:
\begin{align}
    \mathcal{{LOOP}}_{LL}
    &= \sum_{c=\pm} \frac{(-1/2+i\tn)^4}{(\tilde\nu^2 + 1/4)^2} \frac{(\tau_1 \tau_2)^{1+2ic\tn}}{(4\pi)^2} \Gamma^4(-ic\tn) \cdot \frac{4^{-2\ii c\tn} k_s^{3+4ic\tn}}{(4\pi)^{3/2}} \nonumber \\
    &\times \left( f_0(x-1,x+1) +2f_c(x-1/2,x+1) +f_{cc}(x,x+1) \right)
\end{align}
where $x=-\ii c\tn$. Using the identity $\mb q \cdot \mb K=\frac{1}{2}(q^2+K^2-k_s^2)$, we can also express it as:
\begin{align}
    \mathcal{{LOOP}}_{LL}
     &= \sum_{c=\pm} \frac{(-1/2+i\tn)^4}{(\tilde\nu^2 + 1/4)^2} \frac{(\tau_1 \tau_2)^{1+2ic\tn}}{(4\pi)^2} \Gamma^4(-ic\tn) \cdot \frac{4^{-2\ii c\tn} k_s^{3+4ic\tn}}{(4\pi)^{3/2}} \nonumber \\
    &\times \frac{1}{4} \left(f(x-1,x+1) + f(x+1,x-1) + f(x+1,x+1) \right. \nonumber \\
    &\quad \left. +2f(x,x) - 2f(x,x+1) - 2f(x+1,x)  \right)
\end{align}

We denote it as:
\begin{align}
    \mathcal{{LOOP}}_{LL}
    &= \sum_{c=\pm}
    (\tau_1 \tau_2)^{1+2ic\tn} \cdot
    \frac{\Gamma^4(-ic\tn)}{(4\pi)^2}
    \frac{4^{-2\ii c\tn} k_s^{3+4ic\tn}}{(4\pi)^{3/2}} 
    \cdot \mathcal{I}_{LL}(-ic\tn)
\end{align}

And also for the transversal mode:
\begin{align}
    \mathcal{{LOOP}}_{++}
    &= \int \frac{\di \mb{q}^3}{(2\pi)^3} \left[ D^{(+)}(q,\tau_1,\tau_2) e^{(+)i}_{\mb q} e^{(+)j*}_{\mb q} \cdot D^{(+)}(K,\tau_2,\tau_1) e^{(+)j}_{\mb K} e^{(+)i*}_{\mb K} \right]_{\text{NonL}} \nonumber \\
    &= \sum_{c=\pm} \frac{(\tau_1 \tau_2)^{1+2ic\tn}}{(4\pi)^2} \Gamma^4(-ic\tn) \int \frac{\di \mb{q}^3}{(2\pi)^3}\left[  \left( \frac{q}{2} \right)^{2ic\tn} \left( \frac{K}{2} \right)^{2ic\tn} |e^{(+)}_{\mb q} \cdot e^{(+)*}_{\mb K}|^2 \right] 
\end{align}

Also considering the vector structure:
\begin{align}
    \mathcal{{LOOP}}_{++}  &= \sum_{c=\pm} \frac{(\tau_1 \tau_2)^{1+2ic\tn}}{(4\pi)^2} \Gamma^4(-ic\tn) \times \frac{4^{-2\ii c\tn} k_s^{3+4ic\tn}}{(4\pi)^{3/2}} \nonumber \\
    &\times \frac{1}{4} [f_0(x,x) + f_0(x-1,x+1) +f_{cc}(x,x+1) \nonumber \\
    &\quad + 2f_0(x-\fr{1}{2},x+\fr{1}{2}) +2f_c(x-\fr{1}{2},x+1) +2f_c(x,x+\fr{1}{2}) ]\\
    &= \sum_{c=\pm}
    (\tau_1 \tau_2)^{1+2ic\tn} \cdot
    \frac{\Gamma^4(-ic\tn)}{(4\pi)^2}
    \frac{4^{-2\ii c\tn} k_s^{3+4ic\tn}}{(4\pi)^{3/2}} 
    \cdot \mathcal{I}_{++}(-ic\tn)
\end{align}

We note that the $--$ term equals to the $++$ term. 
Then we should deal with the mixing term.
Firstly we consider the $L+$ term:
\begin{align}
    \mathcal{{LOOP}}_{L+} &= \int \frac{\di \mb{q}^3}{(2\pi)^3} \left[ D^{(+)}(q,\tau_1,\tau_2) e^{(+)i}_{\mb q} e^{(+)j*}_{\mb q} \cdot D^{(+)}(K,\tau_2,\tau_1) e^{(+)j}_{\mb K} e^{(+)i*}_{\mb K} \right]_{\text{NonL}} \nonumber \\
    &= \sum_{c=\pm} \frac{(-1/2+i\tn)^2}{(\tilde\nu^2 + 1/4)} \frac{(\tau_1 \tau_2)^{1+2ic\tn}}{(4\pi)^2} \Gamma^4(-ic\tn) \int \frac{\di \mb{q}^3}{(2\pi)^3}\left[  \left( \frac{q}{2} \right)^{2ic\tn} \left( \frac{K}{2} \right)^{2ic\tn} |e^{(L)}_{\mb q} \cdot e^{(+)*}_{\mb K}|^2 \right] 
\end{align}
and:
\begin{align}
    \mathcal{{LOOP}}_{L+} &= \sum_{c=\pm} \frac{(-1/2+i\tn)^2}{(\tilde\nu^2 + 1/4)} \frac{(\tau_1 \tau_2)^{1+2ic\tn}}{(4\pi)^2} \Gamma^4(-ic\tn) \cdot \frac{4^{-2\ii c\tn} k_s^{3+4ic\tn}}{(4\pi)^{3/2}} \nonumber \\
    &\times \frac{1}{2} \left( f_0(x,x+1) - f_{cc}(x,x+1)  \right) \\
    &= \sum_{c=\pm}
    (\tau_1 \tau_2)^{1+2ic\tn} \cdot
    \frac{\Gamma^4(-ic\tn)}{(4\pi)^2}
    \frac{4^{-2\ii c\tn} k_s^{3+4ic\tn}}{(4\pi)^{3/2}} 
    \cdot \mathcal{I}_{L+}(-ic\tn)
\end{align}
and by completing the dot product we note that $L-=L+=-L=+L$.

And for $+-$ term:
\begin{align}
    \mathcal{{LOOP}}_{+-}  &= \sum_{c=\pm} \frac{(\tau_1 \tau_2)^{1+2ic\tn}}{(4\pi)^2} \Gamma^4(-ic\tn) \times \frac{4^{-2\ii c\tn} k_s^{3+4ic\tn}}{(4\pi)^{3/2}} \nonumber \\
    &\times \frac{1}{4} [f_0(x,x) + f_0(x-1,x+1) +f_{cc}(x,x+1) \nonumber \\
    &\quad - 2f_0(x-\fr{1}{2},x+\fr{1}{2}) +2f_c(x-\fr{1}{2},x+1) -2f_c(x,x+\fr{1}{2}) ] \\
    &= \sum_{c=\pm}
    (\tau_1 \tau_2)^{1+2ic\tn} \cdot
    \frac{\Gamma^4(-ic\tn)}{(4\pi)^2}
    \frac{4^{-2\ii c\tn} k_s^{3+4ic\tn}}{(4\pi)^{3/2}} 
    \cdot \mathcal{I}_{+-}(-ic\tn)
\end{align}
and we have $-+=+-$.

So the loop integral:
\begin{align}
    \mathcal{LOOP} = 
    \sum_{\si_1,\si_2=L,\pm}
    \mathcal{{LOOP}}_{\si_1\si_2}  
    =\sum_{c=\pm} (\tau_1\tau_2)^{1+2ic\tn} \cdot \mathcal{B}_c
\end{align}
where the bubble subgraph contribution is:
\begin{align}
    \mathcal{B}_c = \frac{\Gamma^4(-\ii c\tn)}{(4\pi)^2} \frac{4^{-2\ii c \tn} k_s^{3+4ic\tn}}{(4\pi)^{3/2}} 
    \cdot \sum_{\lam_1,\lam_2=L,\pm} 
    \mathcal{I}_{\lam_1\lam_2}(-\ii c\tn)
\end{align}

Then we can consider the left and right subgraph contribution, namely the time integral.
For the left subgraph:
\begin{align}
    \mathcal{T}^{(L)}_c &= (g^2) \sum_a ia \int \frac{\di \tau_1}{|H \tau_1|^2} G_a(k_1;\tau_1) G_a(k_2;\tau_1) (-\tau_1)^{1+2ic\tn} \nonumber \\
    &= (g^2) \frac{H^2}{4k_1^3 k_2^3} \sum_{a=\pm} ia \int_{-\infty}^{0} \di \tau_1 (-\tau_1)^{-1+2\ii c\tn} (1 -iak_{12}\tau_1-k_1k_2\tau_1^2) e^{iak_{12}\tau_1} \nonumber \\
    &= (g^2) \frac{H^2}{4k_1^3 k_2^3} \times k_{12}^{-2ic\tn} \cdot \mathcal{P}^{(1)}_c(\tn)
\end{align}
where the coefficient is:
\begin{align}
    \mathcal{P}^{(1)}_c(\tn) = \frac{2+ic\tn}{2ic\tn} \ii \sinh(\pi c \tn) \Gamma(2+2\ii c\tn)
\end{align}

Also for the right subgraph:
\begin{align}
    \mathcal{T}^{(R)}_c = (g^2) \frac{H^2}{4k_3^3 k_4^3} \times k_{34}^{-2ic\tn} \cdot \mathcal{P}^{(1)}_c(\tn)
\end{align}

So the total bubble graph contribution:
\begin{align}
    \la \delta h^4 \ra^{\text{Bubble}} &= \sum_{c=\pm} \mathcal{T}^{(L)}_c \mathcal{T}^{(R)}_c \mathcal{B}_c \\
    &= (g^2)^2 \frac{H^4 k_s^3}{16 k_1^3 k_2^3 k_3^3 k_4^3} \left( \frac{k_s^2}{4k_{12}k_{34}}\right)^{2\ii\tn} \times C^{\text{Bub}}(\tn) + \text{c.c.}
\end{align}
where the magnitude is:
\begin{align}
    C^{\text{Bub}}(\tn) = 
    [\mathcal{P}^{(1)}_+(\tn)]^2 \cdot
    \frac{\Gamma^4(-\ii \tn)}{(4\pi)^{7/2}} \cdot 
    \sum_{\lam_1,\lam_2=L,\pm} \mathcal{I}_{\lam_1\lam_2}(\tn)|_{c=+}
\end{align}


\section{Triangle Diagram}
For one of the triangle diagrams:
\begin{align}
    \la \delta h^4 \ra^{\text{Tri1}} &= \left( g^2 \right) (g^2h_0)^2 \sum_{a,b,c=\pm} (-iabc) \int \frac{\di \tau_1}{|H \tau_1|^2} \frac{\di \tau_2}{|H \tau_2|^2} \frac{\di \tau_3}{|H \tau_3|^2} G_a(k_1;\tau_1) G_a(k_2;\tau_2) G_c(k_3;\tau_3) G_c(k_4;\tau_3) \nonumber \\
    &\times \int \frac{\di \mb{q}^3}{(2\pi)^3} \left[ D_{ij}(\mb q+\mb k_1,\tau_1,\tau_2) \right]_{ab} \left[ D_{jk}(\mb q+\mb k_s,\tau_2,\tau_3) \right]_{bc} \left[ D_{ki}( \mb{q},\tau_3,\tau_1) \right]_{ca} 
\end{align}

Firstly we deal with the loop integral:
\begin{align}
    \mathcal{LOOP} 
    = \int \frac{\di \mb{q}^3}{(2\pi)^3} 
    \left[ D_{ij}(\mb q+\mb k_1,\tau_1,\tau_2) \right]_{ab} 
    \left[ D_{jk}(\mb q+\mb k_s,\tau_2,\tau_3) \right]_{bc} 
    \left[ D_{ki}( \mb{q},\tau_3,\tau_1) \right]_{ca}
\end{align}

Using the cutting rule, we know that only when $\mb q$ and $\mb k_s$ both become soft, this integral will contribute the nonlocal signal.
\begin{align}
    \mathcal{LOOP} = \int \frac{\di \mb{q}^3}{(2\pi)^3} \left[ D_{ij}(\mb k_1,\tau_1,\tau_2) \right]_{ab} \left[ D_{jk}(\mb q+\mb k_s,\tau_2,\tau_3) \right] \left[ D_{ki}( \mb{q},\tau_3,\tau_1) \right]
\end{align}
Here we use:
\begin{align}
    e_{i,\mathbf{k}_1}^{(L)}
    &= \begin{pmatrix} \sin\theta_1 \\ 0 \\ \cos\theta_1 \end{pmatrix}, \qquad
    e_{i,\mathbf{k}_1}^{(\pm)}
    = \frac{1}{\sqrt{2}}
    \begin{pmatrix}
    \cos\theta_1 \\
    \pm \ii \\
    -\sin\theta_1
    \end{pmatrix}.
\end{align}
and expanding the propagator by helicity basis:
\begin{align}
    \mathcal{LOOP} = \int \frac{\di \mb{q}^3}{(2\pi)^3} &\left[ D^{(L)}(k_1,\tau_1,\tau_2) e^{(L)i}_{\mb k_1} e^{(L)j*}_{\mb k_1} + D^{(\pm)}(k_1,\tau_1,\tau_2) e^{(\pm)i}_{\mb k_1} e^{(\pm)j*}_{\mb k_1} \right]_{ab} \nonumber \\
    \times &\left[ D^{(L)}(K,\tau_2,\tau_3) e^{(L)j}_{\mb K} e^{(L)k*}_{\mb K} + D^{(\pm)}(K,\tau_2,\tau_3) e^{(\pm)j}_{\mb K} e^{(\pm)k*}_{\mb K} \right] \nonumber \\
    \times &\left[ D^{(L)}(q,\tau_3,\tau_1) e^{(L)k}_{\mb q} e^{(L)i*}_{\mb q} + D^{(\pm)}(q,\tau_3,\tau_1) e^{(\pm)k}_{\mb q} e^{(\pm)i*}_{\mb q} \right] \\
    &= \sum_{\lambda,\sigma_1,\sigma_2=L,\pm} 
    \mathcal{LOOP}_{\sigma_1\sigma_2}^{\lambda}
\end{align}
where each porlarization combination term is:
\begin{align}
    \mathcal{LOOP}_{\sigma_1\sigma_2}^{\lambda} &= \int \frac{\di \mb{q}^3}{(2\pi)^3} 
    [D^{(\lambda)}(k_1,\tau_1,\tau_2)]_{ab} 
    D^{(\sigma_1)}(K,\tau_2,\tau_3) 
    D^{(\sigma_2)}(q,\tau_3,\tau_1) \nonumber \\
    &\times e^{(\lambda)i}_{\mb k_1} e^{(\lambda)j*}_{\mb k_1}
    e^{(\sigma_1)j}_{\mb K} e^{(\sigma_1)k*}_{\mb K}
    e^{(\sigma_2)k}_{\mb q} e^{(\sigma_2)i*}_{\mb q} 
\end{align}

Firstly we consider the $\lambda=L$ terms and complete the contraction with the other two propagators. We set
$\mb k_1=k_1(\sin\theta_1,0,\cos\theta_1)$, and still use $\mb k_s=k_s(0,0,1)$ and
$\mb K=\mb q+\mb k_s$. Then the vector structure is:
\begin{align}
    \mathcal{E}_{\si_1\si_2}^L& = e^{(L)i}_{\mb k_1} e^{(L)j*}_{\mb k_1}
    e^{(\si_1)j}_{\mb K}e^{(\si_1)k*}_{\mb K}
    e^{(\si_2)k}_{\mb q}e^{(\si_2)i*}_{\mb q}
    =
    (e^{(L)*}_{\mb k_1}\cdot e^{(\si_1)}_{\mb K})
    (e^{(\si_1)*}_{\mb K}\cdot e^{(\si_2)}_{\mb q})
    (e^{(\si_2)*}_{\mb q}\cdot e^{(L)}_{\mb k_1}) \nonumber \\
    &= A_{L\si_1} B_{\si_1 \si_2} C_{\si_2 L}
\end{align}
where $\si_1,\si_2=L,\pm$. For simplicity we define:
\begin{align}
    A_{LL} &= \frac{q\sin\theta\sin\theta_1\cos\phi+(q\cos\theta+k_s)\cos\theta_1}{K}, \\
    A_{L+} &= \frac{\sin\theta_1\cos\phi(q\cos\theta+k_s)-\ii K\sin\theta_1\sin\phi-q\sin\theta\cos\theta_1}{\sqrt{2}K}, \\
    A_{L-} &= \frac{\sin\theta_1\cos\phi(q\cos\theta+k_s)+\ii K\sin\theta_1\sin\phi-q\sin\theta\cos\theta_1}{\sqrt{2}K}, \\
    C_{LL} &= \sin\theta\sin\theta_1\cos\phi+\cos\theta\cos\theta_1, \\
    C_{+L} &= \frac{\sin\theta_1(\cos\theta\cos\phi+\ii\sin\phi)-\sin\theta\cos\theta_1}{\sqrt{2}}, \\
    C_{-L} &= \frac{\sin\theta_1(\cos\theta\cos\phi-\ii\sin\phi)-\sin\theta\cos\theta_1}{\sqrt{2}},
\end{align}
and:
\begin{align}
    B_{LL} &= \frac{q+k_s\cos\theta}{K}, &
    B_{L+}&=B_{L-} = -\frac{k_s\sin\theta}{\sqrt{2}K}, \\
    B_{+L}&=B_{-L} = \frac{k_s\sin\theta}{\sqrt{2}K}, &
    B_{++}&=B_{--} = \frac{K+q+k_s\cos\theta}{2K}, \\
    B_{+-}&=B_{-+} = \frac{q+k_s\cos\theta-K}{2K}.
\end{align}
So the nine contractions are:
\begin{align}
    \mathcal{E}_{LL}^L &= A_{LL} B_{LL} C_{LL}, &
    \mathcal{E}_{L+}^L &= A_{LL} B_{L+} C_{+L}, &
    \mathcal{E}_{L-}^L &= A_{LL} B_{L-} C_{-L}, \\
    \mathcal{E}_{+L}^L &= A_{L+} B_{+L} C_{LL}, &
    \mathcal{E}_{++}^L &= A_{L+} B_{++} C_{+L}, &
    \mathcal{E}_{+-}^L &= A_{L+} B_{+-} C_{-L}, \\
    \mathcal{E}_{-L}^L &= A_{L-} B_{-L} C_{LL}, &
    \mathcal{E}_{-+}^L &= A_{L-} B_{-+} C_{+L}, &
    \mathcal{E}_{--}^L &= A_{L-} B_{--} C_{-L} .
\end{align}

Then we can complete the loop integral for the terms related to $D^{(L)}(k_1)$ and recognize the nonlocal signal parts: 
\begin{align}
    \mathcal{LOOP}_{\sigma_1\sigma_2}^{\lambda} &= \int \frac{\di \mb{q}^3}{(2\pi)^3} 
    [D^{(\lambda)}(k_1,\tau_1,\tau_2)]_{ab} 
    D^{(\sigma_1)}(K,\tau_2,\tau_3) 
    D^{(\sigma_2)}(q,\tau_3,\tau_1) 
    \cdot \mathcal{E}^L_{\si_1 \si_2} \\
    &= \left[D^{(\lambda)}(k_1,\tau_1,\tau_2)\right]_{ab}
    \cdot \sum_{d=\pm}
    (\tau_1\tau_2)^{1/2+\ii d\tn}(-\tau_3)^{1+2\ii d\tn} \nonumber \\
    &\times \frac{\Gamma[-\ii d \tn]^4}{(4\pi)^2} 
    \be^{\si_1L} \be^{\si_2L}
    \int \frac{\di \mb{q}^3}{(2\pi)^3}
    \left(\frac{q}{2}\right)^{2\ii d \tn}
    \left(\frac{K}{2}\right)^{2\ii d \tn}
    \mathcal{E}^L_{\si_1 \si_2}   
\end{align}
where $\beta^{\si L}$ is a additional factor for the longitudinal mode:
\begin{align}
    \beta^{\si L} = \begin{cases}
    \frac{(-1/2+\ii \tn)^2}{\tn^2+1/4}, & \si=L, \\
    1, & \si=\pm.
    \end{cases}
\end{align}

Since all terms proportional to
$\sin\phi$ or $\cos\phi$ vanish after the azimuthal integral, we only need the $\phi$-averaged angular factors.
For the azimuthal average, it is important to keep the quadratic terms:
\begin{align}
    \overline{\sin\phi}=\overline{\cos\phi}
    =\overline{\sin\phi\cos\phi}=0,\qquad
    \overline{\sin^2\phi}=\overline{\cos^2\phi}=\frac{1}{2}.
\end{align}

We denote $z=\cos\theta$, $s_1=\sin\theta_1$ and $c_1=\cos\theta_1$. The non-vanishing parts are:
\begin{align}
    \overline{\mathcal{E}}_{LL}^L
    &= \frac{q+k_s z}{K^2}
    \left[\frac{1}{2}q(1-z^2)s_1^2+z(qz+k_s)c_1^2\right], \\
    \overline{\mathcal{E}}_{L+}^L
    =\overline{\mathcal{E}}_{L-}^L
    &= -\frac{k_s(1-z^2)}{2K^2}
    \left[\frac{1}{2}qzs_1^2-(qz+k_s)c_1^2\right], \\
    \overline{\mathcal{E}}_{+L}^L
    =\overline{\mathcal{E}}_{-L}^L
    &= \frac{k_s(1-z^2)}{2K^2}
    \left[\frac{1}{2}(qz+k_s)s_1^2-qzc_1^2\right], \\
    \overline{\mathcal{E}}_{++}^L
    =\overline{\mathcal{E}}_{--}^L
    &= \frac{K+q+k_s z}{4K^2}
    \left[\frac{1}{2}\left(z(qz+k_s)+K\right)s_1^2+q(1-z^2)c_1^2\right], \\
    \overline{\mathcal{E}}_{+-}^L
    =\overline{\mathcal{E}}_{-+}^L
    &= \frac{q+k_s z-K}{4K^2}
    \left[\frac{1}{2}\left(z(qz+k_s)-K\right)s_1^2+q(1-z^2)c_1^2\right].
\end{align}


For any monomial generated by these angular factors, we first define the dimensionful integral:
\begin{align}
    \widetilde{\mathcal{J}}_r(m,n)
    &\equiv
    \int \frac{\di\mb q^3}{(2\pi)^3}
    \left(\frac{q}{2}\right)^{2\ii d\tn}
    \left(\frac{K}{2}\right)^{2\ii d\tn}
    q^m K^{-n} z^r \nonumber \\
    &=\frac{4^{-2\ii d\tn}k_s^{3+4\ii d\tn+m-n}}{(4\pi)^{3/2}}
    \mathcal{J}_r(m,n),
\end{align}
where the dimensionless loop seeds are:
\begin{align}
    \mathcal{J}_r(m,n)
    &\equiv f_r\left(x-\frac{m}{2},x+\frac{n}{2}\right),
    \qquad x=-\ii d\tn ,
\end{align}
and $f_1\equiv f_c$, $f_2\equiv f_{cc}$. The $\mathcal{J}_3$ terms below are reduced to the same loop seeds by using
$z=(K^2-q^2-k_s^2)/(2qk_s)$:
\begin{align}
    \mathcal{J}_3(m,n)
    &= \frac{1}{8}
    \left[\mathcal{J}_0(m-3,n-6)-3\mathcal{J}_0(m-1,n-4)
    -3\mathcal{J}_0(m-3,n-4)\right. \nonumber \\
    &\quad
    +3\mathcal{J}_0(m+1,n-2)+6\mathcal{J}_0(m-1,n-2)
    +3\mathcal{J}_0(m-3,n-2) \nonumber \\
    &\quad\left.
    -\mathcal{J}_0(m+3,n)-3\mathcal{J}_0(m+1,n)
    -3\mathcal{J}_0(m-1,n)-\mathcal{J}_0(m-3,n)\right].
\end{align}

We denote the dimensionless coefficients of the loop integrals as $\mathcal{I}_{\si_1\si_2}^{\lam}(-\ii d\tn)$, which are defined by:
\begin{align}
    \mathcal{LOOP}_{\sigma_1\sigma_2}^{\lambda} 
    &= \left[D^{(\lambda)}(k_1,\tau_1,\tau_2)\right]_{ab}
    \cdot \sum_{d=\pm}
    (\tau_1\tau_2)^{1/2+\ii d\tn}(-\tau_3)^{1+2\ii d\tn} 
    \frac{\Gamma^4(-\ii d\tn)}{(4\pi)^{2}}
    \tilde{\mathcal{I}}^{\lam}_{\si_1 \si_2}(-\ii d\tn)
\end{align}
and explicitly:
\begin{align}
    \widetilde{\mathcal{I}}^{\lambda}_{\sigma_1\sigma_2}
    &\equiv \be^{\si_1L} \be^{\si_2L} \int\frac{\di\mb q^3}{(2\pi)^3}
    \left(\frac{q}{2}\right)^{2\ii f\tn}
    \left(\frac{K}{2}\right)^{2\ii f\tn}
    \overline{\mathcal{E}}^{\lambda}_{\sigma_1\sigma_2} \nonumber \\
    &=\frac{4^{-2\ii f\tn}k_s^{3+4\ii f\tn}}{(4\pi)^{3/2}}
    \mathcal{I}^{\lambda}_{\sigma_1\sigma_2},
    \qquad f=\pm.
\end{align}

Then the dimensionless coefficients of the nine loop integrals are:
\begin{align}
    \mathcal{I}_{LL}^{L}
    &= \frac{(-1/2+\ii c\tn)^4}{(\tn^2 + 1/4)^2} \cdot
    \left\{\frac{s_1^2}{2}\left[\mathcal{J}_0(2,2)+\mathcal{J}_1(1,2)
    -\mathcal{J}_2(2,2)-\mathcal{J}_3(1,2)\right] \right. \nonumber \\
    &\quad \left.
    +c_1^2\left[\mathcal{J}_2(2,2)+\mathcal{J}_3(1,2)+\mathcal{J}_1(1,2)+\mathcal{J}_2(0,2)\right]  \right\}, \\
    \mathcal{I}_{L+}^{L}
    =\mathcal{I}_{L-}^{L}
    &= \frac{(-1/2+\ii c\tn)^2}{(\tn^2 + 1/4)} \cdot \left\{ -\frac{s_1^2}{4}\left[\mathcal{J}_1(1,2)-\mathcal{J}_3(1,2)\right] \right. \nonumber \\
    &\quad \left.
    +\frac{c_1^2}{2}\left[\mathcal{J}_1(1,2)+\mathcal{J}_0(0,2)
    -\mathcal{J}_3(1,2)-\mathcal{J}_2(0,2)\right] \right\} , \\
    \mathcal{I}_{+L}^{L}
    =\mathcal{I}_{-L}^{L}
    &= \frac{(-1/2+\ii c\tn)^2}{(\tn^2 + 1/4)} \cdot \left\{ \frac{s_1^2}{4}\left[\mathcal{J}_1(1,2)+\mathcal{J}_0(0,2)
    -\mathcal{J}_3(1,2)-\mathcal{J}_2(0,2)\right] \right. \nonumber \\
    &\quad \left.
    -\frac{c_1^2}{2}\left[\mathcal{J}_1(1,2)-\mathcal{J}_3(1,2)\right] \right\},  \\
    \mathcal{I}_{++}^{L}
    =\mathcal{I}_{--}^{L}
    &= \frac{s_1^2}{8}\left[
    \mathcal{J}_2(1,1)+\mathcal{J}_2(2,2)+\mathcal{J}_3(1,2)
    +\mathcal{J}_1(0,1)+\mathcal{J}_1(1,2)+\mathcal{J}_2(0,2)\right. \nonumber \\
    &\quad\left.
    +\mathcal{J}_0(0,0)+\mathcal{J}_0(1,1)
    +\mathcal{J}_1(0,1)\right] \nonumber \\
    &\quad
    +\frac{c_1^2}{4}\left[\mathcal{J}_0(1,1)-\mathcal{J}_2(1,1)
    +\mathcal{J}_0(2,2)-\mathcal{J}_2(2,2)
    +\mathcal{J}_1(1,2)-\mathcal{J}_3(1,2)\right], \\
    \mathcal{I}_{+-}^{L}
    =\mathcal{I}_{-+}^{L}
    &= \frac{s_1^2}{8}\left[
    \mathcal{J}_2(2,2)+\mathcal{J}_3(1,2)-\mathcal{J}_2(1,1)
    +\mathcal{J}_1(1,2)+\mathcal{J}_2(0,2)-\mathcal{J}_1(0,1) \right. \nonumber \\
    &\quad\left.
    -\mathcal{J}_0(1,1)-\mathcal{J}_1(0,1)
    +\mathcal{J}_0(0,0) \right] \nonumber \\
    &\quad
    +\frac{c_1^2}{4}
    \left[\mathcal{J}_0(2,2)-\mathcal{J}_2(2,2)
    +\mathcal{J}_1(1,2)-\mathcal{J}_3(1,2)
    -\mathcal{J}_0(1,1)+\mathcal{J}_2(1,1)\right].
\end{align}

So this part of the triangle loop integral can be written as:
\begin{align}
    \mathcal{LOOP}^{L} 
    &= \sum_{\si_1,\si_2=L,\pm} 
    \mathcal{LOOP}^{L}_{\si_1\si_2} \nonumber \\
    &= \left[D^{(L)}(k_1,\tau_1,\tau_2)\right]_{ab} \cdot
    \sum_{d=\pm} (\tau_1\tau_2)^{1/2+\ii d\tn}(-\tau_3)^{1+2\ii d\tn} \nonumber \\
    &\times \frac{4^{-2\ii d\tn} k_s^{3+4\ii d\tn} \Gamma^4(-\ii d\tn)}{(4\pi)^{7/2}}
    \mathcal{I}^{L},
\end{align}
where we define $\mathcal{I}^L$ is the sum of the nine contractions:
\begin{align}
    \mathcal{I}^{L}
    &\equiv \sum_{\si_1,\si_2=L,\pm} \mathcal{I}^{L}_{\si_1\si_2}
    = (\mathcal{I}^{L}_{LL}
    +2\mathcal{I}^{L}_{L+}
    +2\mathcal{I}^{L}_{+L}
    +2\mathcal{I}^{L}_{++}
    +2\mathcal{I}^{L}_{+-})
\end{align}

Equivalently:
\begin{align}
    \mathcal{LOOP}^{L}
    =
    \left[D^{(L)}(k_1,\tau_1,\tau_2)\right]_{ab} \cdot
    \sum_{d=\pm}
    (\tau_1\tau_2)^{1/2+\ii d\tn}
    (-\tau_3)^{1+2\ii d\tn}
    \mathcal{B}^L_d,
\end{align}
with:
\begin{align}
    \mathcal{B}^L_d
    =\frac{4^{-2\ii d\tn}k_s^{3+4\ii d\tn} \Gamma^4(-\ii d\tn)}{(4\pi)^{7/2}}
    \mathcal{I}^{(L)}.
\end{align}

Now we repeat the same calculation for $\lambda=\pm$ terms. 
The vector contractions are:
\begin{align}
    \mathcal{E}^{(\sigma)}_{\si_1\si_2}
    &=
    e^{(\sigma)i}_{\mb k_1} e^{(\sigma)j*}_{\mb k_1}
    e^{(\si_1)j}_{\mb K}e^{(\si_1)k*}_{\mb K}
    e^{(\si_2)k}_{\mb q}e^{(\si_2)i*}_{\mb q} \nonumber \\
    &=
    (e^{(\sigma)*}_{\mb k_1}\cdot e^{(\si_1)}_{\mb K})
    (e^{(\si_1)*}_{\mb K}\cdot e^{(\si_2)}_{\mb q})
    (e^{(\si_2)*}_{\mb q}\cdot e^{(\sigma)}_{\mb k_1}),
    \qquad \sigma=\pm .
\end{align}
For the first and third factors we define:
\begin{align}
    A_{+L}
    &=\frac{q\sin\theta\cos\phi\cos\theta_1-\ii q\sin\theta\sin\phi-(q\cos\theta+k_s)\sin\theta_1}{\sqrt{2}K}, \\
    A_{++}
    &=\frac{(q\cos\theta+k_s)\cos\theta_1\cos\phi
    - \ii K \sin\phi\cos\theta_1
    - \ii \sin\phi (q\cos\theta+k_s)
    + K \cos\phi
    + q \sin\theta\sin\theta_1}
    {2K}, \nonumber \\
    A_{+-}
    &=\frac{(q\cos\theta+k_s)\cos\theta_1\cos\phi
    + \ii K \sin\phi\cos\theta_1
    - \ii \sin\phi (q\cos\theta+k_s)
    - K \cos\phi
    + q \sin\theta\sin\theta_1}
    {2K}, \nonumber \\
    C_{L+}
    &=\frac{\sin\theta\cos\phi\cos\theta_1
    +\ii\sin\theta\sin\phi
    -\cos\theta\sin\theta_1}
    {\sqrt{2}}, \nonumber \\
    C_{++}
    &=\frac{\cos\theta\cos\theta_1\cos\phi
    + \ii \sin\phi \cos\theta_1
    + \ii \cos\theta\sin\phi
    + \cos\phi
    + \sin\theta\sin\theta_1}
    {2}, \nonumber \\
    C_{-+}
    &=\frac{\cos\theta\cos\theta_1\cos\phi
    - \ii \sin\phi \cos\theta_1
    + \ii \cos\theta\sin\phi
    - \cos\phi
    + \sin\theta\sin\theta_1}
    {2}. 
\end{align}
The $D^{(-)}(k_1)$ factors are obtained by complex conjugating the external polarization:
\begin{align}
    A_{-L}
    &=\frac{q\sin\theta\cos\phi\cos\theta_1
    +\ii q\sin\theta\sin\phi
    -(q\cos\theta+k_s)\sin\theta_1}{\sqrt{2}K}, \nonumber \\
    A_{-+}
    &=\frac{(q\cos\theta+k_s)\cos\theta_1\cos\phi
    - \ii K \sin\phi\cos\theta_1
    + \ii \sin\phi (q\cos\theta+k_s)
    - K \cos\phi
    + q \sin\theta\sin\theta_1}
    {2K}, \nonumber \\
    A_{--}
    &=\frac{(q\cos\theta+k_s)\cos\theta_1\cos\phi
    + \ii K \sin\phi\cos\theta_1
    + \ii \sin\phi (q\cos\theta+k_s)
    + K \cos\phi
    + q \sin\theta\sin\theta_1}
    {2K}, \nonumber \\
    C_{L-}
    &=\frac{\sin\theta\cos\phi\cos\theta_1
    -\ii\sin\theta\sin\phi
    -\cos\theta\sin\theta_1}
    {\sqrt{2}}, \nonumber \\
    C_{+-}
    &=\frac{\cos\theta\cos\theta_1\cos\phi
    + \ii \sin\phi \cos\theta_1
    - \ii \cos\theta\sin\phi
    - \cos\phi
    + \sin\theta\sin\theta_1}
    {2}, \nonumber \\
    C_{--}
    &=\frac{\cos\theta\cos\theta_1\cos\phi
    - \ii \sin\phi \cos\theta_1
    - \ii \cos\theta\sin\phi
    + \cos\phi
    + \sin\theta\sin\theta_1}
    {2}.
\end{align}
The middle factor $B_{\si_1\si_2}$ is the same as before. Therefore the nine contractions are:
\begin{align}
    \mathcal{E}^{\sigma}_{LL} &= A_{\si L} B_{LL} C_{L\si}, &
    \mathcal{E}^{\si}_{L+} &= A_{\si L} B_{L+} C_{+\si}, &
    \mathcal{E}^{\si}_{L-} &= A_{\si L} B_{L-} C_{-\si}, \\
    \mathcal{E}^{\si}_{+L} &= A_{\si +} B_{+L} C_{L\si}, &
    \mathcal{E}^{\si}_{++} &= A_{\si +} B_{++} C_{+\si}, &
    \mathcal{E}^{\si}_{+-} &= A_{\si +} B_{+-} C_{-\si}, \\
    \mathcal{E}^{\si}_{-L} &= A_{\si -} B_{-L} C_{L\si}, &
    \mathcal{E}^{\si}_{-+} &= A_{\si -} B_{-+} C_{+\si}, &
    \mathcal{E}^{\si}_{--} &= A_{\si -} B_{--} C_{-\si} .
\end{align}

After the azimuthal average, we can get the non-vanishing parts of the nine contractions:
\begin{align}
    \bar{\mathcal{E}}^+_{LL} 
    &= \frac{q+k_s z}{2K^2}
    \left[
        \frac{q(1+c_1^2)}{2} (1-z^2) 
        + (qz+k_s)z s_1^2
    \right] \nonumber \\
    \bar{\mathcal{E}}^+_{L+}
    &= \left(-\frac{k_s(1-z^2)}{4K^2}\right)
    \left[
        qc_1 + qz\frac{1+c_1^2}{2} -(qz+k_s)s_1^2
    \right] \nonumber \\
    \bar{\mathcal{E}}^+_{L-}
    &= \left(-\frac{k_s(1-z^2)}{4K^2}\right)
    \left[
        -qc_1 + qz\frac{1+c_1^2}{2} -(qz+k_s)s_1^2
    \right] \nonumber \\
    \bar{\mathcal{E}}^+_{+L}
    &= \left(\frac{k_s(1-z^2)}{4K^2}\right)
    \left[
        Kc_1 + (qz+k_s)\frac{1+c_1^2}{2} -qz s_1^2
    \right] \nonumber \\
    \bar{\mathcal{E}}^+_{-L}
    &= \left(\frac{k_s(1-z^2)}{4K^2}\right)
    \left[
        -Kc_1 + (qz+k_s)\frac{1+c_1^2}{2} -qz s_1^2
    \right] \nonumber \\
    \bar{\mathcal{E}}^+_{++}
    &= \left(\frac{K+q+k_sz}{8K^2}\right)
    \left[
        (qz+k_s)\frac{c_1^2z+2c_1+z}{2} +K\frac{c_1^2+2c_1+1}{2} + qs_1^2(1-z^2)
    \right] \nonumber \\
    \bar{\mathcal{E}}^+_{--}
    &= \left(\frac{K+q+k_sz}{8K^2}\right)
    \left[
        (qz+k_s)\frac{c_1^2z-2c_1+z}{2} +K\frac{c_1^2-2c_1+1}{2} + qs_1^2(1-z^2)
    \right] \nonumber \\
    \bar{\mathcal{E}}^+_{+-}
    &= \left(\frac{q+k_sz-K}{8K^2}\right)
    \left[
        (qz+k_s)\frac{c_1^2z-2c_1+z}{2} -K\frac{c_1^2-2zc_1+1}{2} + qs_1^2(1-z^2)
    \right] \nonumber \\
    \bar{\mathcal{E}}^+_{-+}
    &= \left(\frac{q+k_sz-K}{8K^2}\right)
    \left[
        (qz+k_s)\frac{c_1^2z+2c_1+z}{2} -K\frac{c_1^2+2zc_1+1}{2} + qs_1^2(1-z^2)
    \right]
\end{align}

Applying the same loop-seed map as above, we get the dimensionless coefficient:
\begin{align}
    {\mathcal{I}}^+_{LL} 
    &= \frac{(-1/2 + \ii d \tn)^4}{(\tn^2+1/4)^2} 
    \cdot \left\{ \frac{1+c_1^2}{4} [ \mathcal{J}_0(2,2) -\mathcal{J}_2(2,2) +\mathcal{J}_1(1,2) -\mathcal{J}_3(1,2) ] \right. \nonumber \\
    &\qquad \left. +\frac{s_1^2}{2} [\mathcal{J}_2(2,2) +\mathcal{J}_1(1,2) +\mathcal{J}_3(1,2) +\mathcal{J}_2(0,2) ] \right\} \nonumber \\
    {\mathcal{I}}^+_{L+} + {\mathcal{I}}^+_{L-} 
    + {\mathcal{I}}^+_{+L} + {\mathcal{I}}^+_{-L} 
    &= \frac{(-1/2 + \ii d \tn)^2}{(\tn^2+1/4)}
    \frac{2+s_1^2}{4} [\mathcal{J}_0(0,2)-\mathcal{J}_2(0,2)] \nonumber \\
    {\mathcal{I}}^+_{++} + {\mathcal{I}}^+_{--} 
    + {\mathcal{I}}^+_{+-} + {\mathcal{I}}^+_{-+} 
    &= \frac{1+c_1^2}{4} [\mathcal{J}_2(2,2) +\mathcal{J}_1(1,2)
    +\J_3(1,2) +\J_2(0,2) +\J_0(0,0)] \nonumber \\
    &\qquad + \frac{s_1^2}{2} [\J_0(2,2) -\J_2(2,2) +\J_1(1,2) -\J_3(1,2)]     
\end{align}

So the sum of these 9 contraction:
\begin{align}
    \mathcal{I}^{+}
    &\equiv \sum_{\si_1,\si_2=L,\pm} \mathcal{I}^{+}_{\si_1\si_2}
\end{align}

Also for the $D^{(-)}(k_1)$ factors, the polarization contractions are:
\begin{align}
    \bar{\mathcal{E}}^{-}_{LL} &= \bar{\mathcal{E}}^{+}_{LL}, &
    \bar{\mathcal{E}}^{-}_{L+} &= \bar{\mathcal{E}}^{+}_{L-}, &
    \bar{\mathcal{E}}^{-}_{L-} &= \bar{\mathcal{E}}^{+}_{L+}, \\
    \bar{\mathcal{E}}^{-}_{+L} &= \bar{\mathcal{E}}^{+}_{-L}, &
    \bar{\mathcal{E}}^{-}_{-L} &= \bar{\mathcal{E}}^{+}_{+L}, &
    \bar{\mathcal{E}}^{-}_{++} &= \bar{\mathcal{E}}^{+}_{--}, \\
    \bar{\mathcal{E}}^{-}_{--} &= \bar{\mathcal{E}}^{+}_{++}, &
    \bar{\mathcal{E}}^{-}_{+-} &= \bar{\mathcal{E}}^{+}_{-+}, &
    \bar{\mathcal{E}}^{-}_{-+} &= \bar{\mathcal{E}}^{+}_{+-} .
\end{align}

So the sum of these 9 contraction:
\begin{align}
    \mathcal{I}^{-}
    &\equiv \sum_{\si_1,\si_2=L,\pm} \mathcal{I}^{-}_{\si_1\si_2}
    =\mc I^+
\end{align}

So the transverse part of the triangle loop integral is:
\begin{align}
    \mathcal{LOOP}^{\si}
    =
    \left[D^{(\sigma)}(k_1,\tau_1,\tau_2)\right]_{ab} \cdot
    \sum_{d=\pm}(\tau_1\tau_2)^{1/2+\ii d\tn} (-\tau_3)^{1+2\ii d\tn}
    \mathcal{B}^{\si}_d,
\end{align}
with the bubble subgraph contribution:
\begin{align}
    \mathcal{B}^{\si}_d
    =
    \frac{\Gamma^4(-\ii d\tn)}{(4\pi)^{7/2}}
    4^{-2\ii d\tn}k_s^{3+4\ii d\tn}
    \mathcal{I}^{\si},
    \qquad \sigma=\pm .
\end{align}

So the total loop integral becomes:
\begin{align}
    \mathcal{LOOP}
    =
    \sum_{\lambda=T,\pm}
    \left[D^{(\lambda)}(k_1,\tau_1,\tau_2)\right]_{ab}
    \sum_{d=\pm}(\tau_1\tau_2)^{1/2+\ii d\tn} (-\tau_3)^{1+2\ii d\tn}
    \mathcal{B}^{\lambda}_d
\end{align}

Then we also deal with the left and right subgraph contribution:
\begin{align}
    \mathcal{T}^{(L),\lambda}_d &= (g^2h_0)^2 \sum_{a,b=\pm} (-ab) \int \frac{\di \tau_1}{|H \tau_1|^2} \frac{\di \tau_2}{|H \tau_2|^2} G_a(k_1;\tau_1) G_b(k_2;\tau_2) \nonumber \\
    &\times \left[D^{(\lambda)}(k_1,\tau_1,\tau_2)\right]_{ab} (\tau_1\tau_2)^{1/2+\ii d\tn} \\
    &= (g^2h_0)^2  
    \frac{H^2}{4k_1^3 k_2^3}
    \sum_{a,b=\pm} (-ab)
    \int \di \tau_1 \di \tau_2 
    (1-iak_1\tau_1)(1-ibk_2\tau_2)
    \nonumber \\
    &\times \left[D^{(\lambda)}(k_1,\tau_1,\tau_2)\right]_{ab} 
    (-\tau_1)^{-3/2+\ii d\tn}
    (-\tau_2)^{-3/2+\ii d\tn} 
    e^{\ii ak_1\tau_1 + \ii bk_2\tau_2}
\end{align}

For $\lambda=L$, this integral is given by the tree seed in \cite{qin_helical_2023}:
\begin{align}
    \mathcal{T}^{(L),L}_d 
    &= \left( g^2h_0\right)^2 
    \frac{k_{12}^{-2id\tn}}{4k_1^3 k_2^3} 
    \cdot \mathcal{P}_d^{(2),L}(\tn)
\end{align}
where the coefficient is:
\begin{align}
    \mathcal{P}_d^{(2),L}(\tn) 
    = 2^{2id\tn} \cdot \sum_{a,b=\pm} 
    &\left[ +{\mathcal{K}}_{ab}^{-\fr32+id\tn,-\fr32+id\tn}  
    - ia {\mathcal{K}}_{ab}^{-\fr12+id\tn,-\fr32+id\tn} \right. \nonumber \\
    &\left. -ib {\mathcal{K}}_{ab}^{-\fr32+id\tn,-\fr12+id\tn} 
    - ab {\mathcal{K}}_{ab}^{-\fr12+id\tn,-\fr12+id\tn} \right]
\end{align}
and here $\mathcal{K}_{ab}^{p_1p_2} = \mathcal{I}^{(L)p_1p_2}_{ab}(1,1)$ are vector seeds at the double folded limit in \cite{qin_helical_2023}. The opposite-sign seeds:
\begin{align}
    \mathcal{K}_{\pm\mp}^{p_1p_2} = \frac{e^{\mp i\pi(p_1-p_2)/2}}{4\pi(\tilde{\nu}^2 + \frac{1}{4})} \left[ \mathbf{H}_{\tilde{\nu}}^{p_1}(1) \left(\frac{1}{2}\right)^{i\tilde{\nu}} + \mathbf{H}_{-\tilde{\nu}}^{p_1}(1) \left(\frac{1}{2}\right)^{-i\tilde{\nu}} \right] \nonumber \\
    \times \left[ \mathbf{H}_{\tilde{\nu}}^{p_2}(1) \left(\frac{1}{2}\right)^{i\tilde{\nu}} + \mathbf{H}_{-\tilde{\nu}}^{p_2}(1) \left(\frac{1}{2}\right)^{-i\tilde{\nu}} \right]
\end{align}
where we define:
\begin{align}
    \mathbf{H}_{\tilde{\nu}}^p(1) \equiv \Bigg\{ \left( \frac{1}{2} - i\tilde{\nu} \right) \Gamma\left[ -i\tilde{\nu}, p + \frac{3}{2} + i\tilde{\nu} \right] {}_2F_1\left[ \begin{matrix} \frac{p}{2} + \frac{3}{4} + \frac{i\tilde{\nu}}{2}, & \frac{p}{2} + \frac{5}{4} + \frac{i\tilde{\nu}}{2} \\ 1 + i\tilde{\nu} \end{matrix} \middle| 1 \right] \nonumber \\
    + \frac{1}{2} \Gamma\left[ -1 - i\tilde{\nu}, p + \frac{7}{2} + i\tilde{\nu} \right] {}_2F_1\left[ \begin{matrix} \frac{p}{2} + \frac{7}{4} + \frac{i\tilde{\nu}}{2}, & \frac{p}{2} + \frac{9}{4} + \frac{i\tilde{\nu}}{2} \\ 2 + i\tilde{\nu} \end{matrix} \middle| 1 \right] \Bigg\}
\end{align}
And the same-sign seeds:
\begin{align}
    \mathcal{K}_{\pm\pm}^{p_1p_2}
    &= \frac{-i e^{\mp i\pi(p_1+p_2)/2}}{4\pi(\tilde{\nu}^2 + \frac{1}{4})} 
    \left[ e^{\pi\tilde{\nu}} \mathbf{H}_{\tilde{\nu}}^{p_1}(1) \left(\frac{1}{2}\right)^{i\tilde{\nu}} 
    + e^{-\pi\tilde{\nu}} \mathbf{H}_{-\tilde{\nu}}^{p_1}(1) \left(\frac{1}{2}\right)^{-i\tilde{\nu}} \right]  \left[ \mathbf{H}_{\tilde{\nu}}^{p_2}(1) \left(\frac{1}{2}\right)^{i\tilde{\nu}} + \mathbf{H}_{-\tilde{\nu}}^{p_2}(1) \left(\frac{1}{2}\right)^{-i\tilde{\nu}} \right] \nonumber \\
    &\quad + \frac{e^{\mp i\pi(p_1+p_2)/2}}{2(\tilde{\nu}^2 + \frac{1}{4})}  
    \sum_{n_1,n_2=0}^{\infty} \frac{(-1)^{n_{12}}}{n_1! n_2!} \left( \frac{1}{2} \right)^{2n_{12}} \frac{i}{\tilde{\nu}} \left( 2n_1 - \frac{1}{2} + i\tilde{\nu} \right) \left( 2n_2 - \frac{1}{2} - i\tilde{\nu} \right) (-i\tilde{\nu})_{-n_1} (i\tilde{\nu})_{-n_2} \nonumber \\
    &\quad \times \mathcal{F} \left[ \begin{matrix} 2n_2 + p_2 + \frac{3}{2} - i\tilde{\nu}, & 2n_{12} + p_1 + p_2 + 3 \\ 2n_2 + p_2 + \frac{5}{2} - i\tilde{\nu} \end{matrix} \middle| -1 \right] + (\tilde{\nu} \rightarrow -\tilde{\nu})
\end{align}

But in this way, the double folded limit is hard to get.
So we use the other method, namely acting differential operator on the scalar seed\cite{qin_closed-form_2023}.
Then the coefficient becomes two parts, where the first part contains delta function contribution:
\begin{align}
    \text{Part 1 of}\;\; \mathcal{P}_d^{(2),L}(\tn) = \mathcal{P}_d^{(1)}(\tn) \cdot \frac{1}{\tn^2+1/4}
\end{align}

And the second part comes from the scalar seeds:
\begin{align}
    \text{Part 2 of}\;\; \mathcal{P}_d^{(2),L}(\tn) &= 
    \frac{2^{2id\tn}}{\tn^2+1/4} \sum_{a,b=\pm} 
    \left\{ (\fr12+id\tn)^2 \left[ 
    \mc M_{ab}^{p_0,p_0}
    +(ia) \mc M_{ab}^{p_0+1,p_0}
    +(ib) \mc M_{ab}^{p_0,p_0+1}
    +(ab) \mc M_{ab}^{p_0+1,p_0+1}
    \right] \right. \nonumber \\
    &+ (\fr12+id\tn) \left[
    \mc M_{ab}^{p_0+2,p_0}
    + (ib) \mc M_{ab}^{p_0+2,p_0+1}
    + \mc M_{ab}^{p_0,p_0+2}
    + (ia) \mc M_{ab}^{p_0+1,p_0+2}
    \right] \nonumber \\
    &+ \left. \mc M_{ab}^{p_0+2,p_0+2} \right\}
\end{align}
where $p_0=-\fr52+\ii d\tn$, and  $\mathcal{M}_{ab}^{p_1p_2} = \mathcal{I}^{p_1p_2}_{ab}(1,1)$ are the scalar seeds in \cite{qin_closed-form_2023}.
The opposite-sign seeds:
\begin{align}
    \mc M^{p_1 p_2}_{\pm\mp} = 
    \frac{e^{\mp \mathrm{i} \bar{p}_{12} \pi/2}}
    {2^{5+p_{12}}} 
    \Gamma \left[ \begin{matrix} 
        \frac{5}{2} + p_1 - \mathrm{i}\tilde{\nu}, 
        & \frac{5}{2} + p_1 + \mathrm{i}\tilde{\nu}, 
        & \frac{5}{2} + p_2 - \mathrm{i}\tilde{\nu}, 
        & \frac{5}{2} + p_2 + \mathrm{i}\tilde{\nu} \\ 
        3 + p_1, 
        & 3 + p_2 
    \end{matrix} \right]
\end{align}

And the same-sign seeds:
\begin{align}
    \mathcal{M}^{p_1 p_2}_{\pm\pm} 
    &= \frac{\pm \mathrm{i} e^{\mp \mathrm{i} p_{12} \pi/2} e^{-\pi \widetilde{\nu}}}{2^{5+p_{12}}} 
    \Gamma \begin{bmatrix} 
    \frac{5}{2} + p_1 - \mathrm{i}\widetilde{\nu}, & \frac{5}{2} + p_1 + \mathrm{i}\widetilde{\nu}, & \frac{5}{2} + p_2 - \mathrm{i}\widetilde{\nu}, & \frac{5}{2} + p_2 + \mathrm{i}\widetilde{\nu} \\ 
    3 + p_1, & 3 + p_2 
    \end{bmatrix} \nonumber \\
    &\quad - \frac{e^{\mp \mathrm{i} p_{12} \pi/2}}{2^{5+p_{12}}} 
    \Gamma \begin{bmatrix} 
    5 + p_{12}, & \frac{5}{2} + p_1 \pm \mathrm{i}\widetilde{\nu}, & \frac{5}{2} + p_2 \pm \mathrm{i}\widetilde{\nu} 
    \end{bmatrix} \nonumber \\
    &\quad \cdot {}_3\widetilde{F}_2 \left[ \begin{matrix} 
    5 + p_{12}, & \frac{1}{2} \pm \mathrm{i}\widetilde{\nu}, & 1 \\ 
    \frac{7}{2} + p_1 \pm \mathrm{i}\widetilde{\nu}, & \frac{7}{2} + p_2 \pm \mathrm{i}\widetilde{\nu} 
    \end{matrix} \middle| 1 \right].
\end{align}

For $\lambda=\pm$, this integral can also be given by the tree seed in \cite{qin_closed-form_2023}:
\begin{align}
    \mathcal{T}^{(L),\pm}_d 
    &= (g^2h_0)^2  
    \frac{H^2}{4k_1^3 k_2^3}
    \sum_{a,b=\pm} (-ab)
    \int \di \tau_1 \di \tau_2 
    (1-iak_1\tau_1)(1-ibk_2\tau_2)
    \nonumber \\
    &\times \left[D^{(\text{scalar})}(k_1,\tau_1,\tau_2)\right]_{ab} 
    (-\tau_1)^{-5/2+\ii d\tn}
    (-\tau_2)^{-5/2+\ii d\tn} 
    e^{\ii ak_1\tau_1 + \ii bk_2\tau_2} \\
    &= \left( g^2h_0\right)^2 
    \frac{k_{12}^{-2id\tn}}{4k_1^3 k_2^3} 
    \cdot \mathcal{P}_d^{(2),\pm}(\tn)
\end{align}
where the coefficient is:
\begin{align}
    \mathcal{P}_d^{(2),\pm}(\tn) 
    = 2^{2id\tn} \cdot \sum_{a,b=\pm} 
    &\left[ +{\mathcal{M}}_{ab}^{p_0,p_0} 
    - ia {\mathcal{M}}_{ab}^{p_0+1,p_0}  
    - ib {\mathcal{M}}_{ab}^{p_0,p_0+1}  
    - ab {\mathcal{M}}_{ab}^{p_0+1,p_0+1} \right]
\end{align}

And for the right subgraph:
\begin{align}
    \mathcal{T}^{(R)}_d &= (g^2) \sum_c ic \int \frac{\di \tau_3}{|H \tau_3|^2} G_c(k_3;\tau_1) G_c(k_4;\tau_3) (-\tau_3)^{1+2\ii d\tn} \nonumber \\
    &= (g^2) \frac{H^2}{4k_3^3 k_4^3} \sum_{c=\pm} ic \int_{-\infty}^{0} \di \tau_3 (-\tau_3)^{-1+2\ii c\tn} (1 -ick_{34}\tau_3-k_3k_4\tau_3^2) e^{ick_{34}\tau_3} \nonumber \\
    &= (g^2) \frac{H^2}{4k_3^3 k_4^3} \times k_{34}^{-2\ii d\tn} \cdot \mathcal{P}^{(1)}_d(\tn)
\end{align}
where the coefficient is:
\begin{align}
    \mathcal{P}^{(1)}_d(\tn) = \frac{2+\ii d\tn}{2\ii d\tn} \ii\sinh(\pi d \tn) \Gamma(2+2\ii d\tn)
\end{align}

So one of the triangle diagram:
\begin{align}
    \la \delta h^4 \ra^{\text{Tri1}} 
    &= \sum_{d=\pm} \sum_{\lambda=L,\pm} 
    \mc T^{(L),\lambda}_d  
    \cdot \mc B^{\lam}_d 
    \cdot \mc T^{(R)}_d \\
    &= (g^2)(g^2h_0)^2 
    \frac{H^2 k_s^3}{16 k_1^3 k_2^3 k_3^3 k_4^3} 
    \left( \frac{k_s^2}{4k_{12}k_{34}}\right)^{2\ii\tn} 
    \times C^{\text{Tri}}(\tn,\theta_1) + \text{c.c.}
\end{align}
where the magnitude is:
\begin{align}
    C^{\text{Tri}}(\tn,\theta_1) = 
    \sum_{\lam=L,\pm}
    [\mathcal{P}^{(1)}_+(\tn) \mc P^{(2),\lambda}_+(\tn) ]
    \cdot \frac{\Gamma^4(-\ii \tn)}{(4\pi)^{7/2}} \cdot 
    \mathcal{I}^{\lambda}(\tn,\theta_1)|_{d=+}
\end{align}

Also for another triangle:
\begin{align}
    \la \delta h^4 \ra^{\text{Tri2}} 
    = \la \delta h^4 \ra^{\text{Tri1}}|_{k_1 \longleftrightarrow k_3, k_2 \longleftrightarrow k_4, \theta_1 \longleftrightarrow \theta_3, } 
\end{align}

So the total triangle diagram contributions:
\begin{align}
    \la \delta h^4 \ra^{\text{Tri}} 
    &= \la \delta h^4 \ra^{\text{Tri1}} 
    +\la \delta h^4 \ra^{\text{Tri2}} \nonumber \\
    &= (g^2)(g^2h_0)^2 
    \frac{H^2 k_s^3}{16 k_1^3 k_2^3 k_3^3 k_4^3} 
    \left( \frac{k_s^2}{4k_{12}k_{34}}\right)^{2\ii\tn} 
    \times (C^{\text{Tri}}(\tn,\theta_1) + C^{\text{Tri}}(\tn,\theta_3)) 
    + \text{c.c.}
\end{align}


\section{Box Diagram}
For the box diagrams:
\begin{align}
    \la \delta h^4 \ra^{\text{Box}} &= (ig^2h_0)^4 \sum_{a,b,c,d=\pm} (abcd) 
    \int \frac{\di \tau_1}{|H \tau_1|^2} \frac{\di \tau_2}{|H \tau_2|^2} 
    \frac{\di \tau_3}{|H \tau_3|^2} \frac{\di \tau_4}{|H \tau_4|^2} \nonumber \\
    &\times G_a(k_1;\tau_1) G_a(k_2;\tau_2) G_c(k_3;\tau_3) G_d(k_4;\tau_4) \nonumber \\
    &\times \int \frac{\di \mb{q}^3}{(2\pi)^3} 
    \left[ D_{ij}(\mb q+\mb k_1,\tau_1,\tau_2) \right]_{ab} 
    \left[ D_{jk}(\mb q+\mb k_s,\tau_2,\tau_3) \right]_{bc}  \nonumber \\
    &\times \left[ D_{kl}( \mb{q}+\mb k_3,\tau_3,\tau_4) \right]_{cd} 
    \left[ D_{li}( \mb{q},\tau_4,\tau_1) \right]_{da}
\end{align}

Also we deal with the loop integral:
\begin{align}
    \mathcal{LOOP} &= 
    \int \frac{\di \mb{q}^3}{(2\pi)^3} 
    \left[ D_{ij}(\mb q+\mb k_1,\tau_1,\tau_2) \right]_{ab} 
    \left[ D_{jk}(\mb q+\mb k_s,\tau_2,\tau_3) \right]_{bc}  \nonumber \\
    &\times \left[ D_{kl}( \mb{q}+\mb k_3,\tau_3,\tau_4) \right]_{cd} 
    \left[ D_{li}( \mb{q},\tau_4,\tau_1) \right]_{da}
\end{align}

Using the cutting rule:
\begin{align}
    \mathcal{LOOP} &= 
    \int \frac{\di \mb{q}^3}{(2\pi)^3} 
    \left[ D_{ij}(\mb k_1,\tau_1,\tau_2) \right]_{ab} 
    \left[ D_{jk}(\mb q+\mb k_s,\tau_2,\tau_3) \right] \nonumber \\
    &\times \left[ D_{kl}( \mb{k}_3,\tau_3,\tau_4) \right]_{cd}
    \left[ D_{li}( \mb{q},\tau_4,\tau_1) \right]
\end{align}
Expanding the four propagators in the helicity basis, the loop integral becomes:
\begin{align}
    \mathcal{LOOP}
    =\int\frac{\di\mb q^3}{(2\pi)^3}
    &\left[D^{(L)}(k_1,\tau_1,\tau_2)
    e^{(L)i}_{\mb k_1}e^{(L)j*}_{\mb k_1}
    +D^{(\pm)}(k_1,\tau_1,\tau_2)
    e^{(\pm)i}_{\mb k_1}e^{(\pm)j*}_{\mb k_1}\right]_{ab} \nonumber \\
    &\times\left[D^{(L)}(K,\tau_2,\tau_3)
    e^{(L)j}_{\mb K}e^{(L)k*}_{\mb K}
    +D^{(\pm)}(K,\tau_2,\tau_3)
    e^{(\pm)j}_{\mb K}e^{(\pm)k*}_{\mb K}\right] \nonumber \\
    &\times\left[D^{(L)}(k_3,\tau_3,\tau_4)
    e^{(L)k}_{\mb k_3}e^{(L)l*}_{\mb k_3}
    +D^{(\pm)}(k_3,\tau_3,\tau_4)
    e^{(\pm)k}_{\mb k_3}e^{(\pm)l*}_{\mb k_3}\right]_{cd} \nonumber \\
    &\times\left[D^{(L)}(q,\tau_4,\tau_1)
    e^{(L)l}_{\mb q}e^{(L)i*}_{\mb q}
    +D^{(\pm)}(q,\tau_4,\tau_1)
    e^{(\pm)l}_{\mb q}e^{(\pm)i*}_{\mb q}\right],
\end{align}
where $\mb K=\mb q+\mb k_s$, and each repeated $\pm$ label is summed over the two transverse helicities.

To organize the nine external-helicity combinations, we choose the same convention for $\mb k_1$ as in
the triangle diagram and set
\begin{align}
    \mb k_1&=k_1(\sin\theta_1,0,\cos\theta_1), \\
    \mb k_3&=k_3(\sin\theta_3\cos\phi_3,\sin\theta_3\sin\phi_3,\cos\theta_3).
\end{align}
Let $\lambda_1,\lambda_2=L,\pm$ label the polarizations of the hard propagators with momenta
$\mb k_1$ and $\mb k_3$, respectively, and let $\sigma_1,\sigma_2=L,\pm$ label the two soft propagators.
Keeping the latter two helicities explicit, as in the Triangle calculation, define
\begin{align}
    \mathcal{E}^{\lambda_1\lambda_2}_{\sigma_1\sigma_2}
    &\equiv
    e_{\mb k_1}^{(\lambda_1)i}e_{\mb k_1}^{(\lambda_1)j*}
    e_{\mb K}^{(\sigma_1)j}e_{\mb K}^{(\sigma_1)k*}
    e_{\mb k_3}^{(\lambda_2)k}e_{\mb k_3}^{(\lambda_2)l*}
    e_{\mb q}^{(\sigma_2)l}e_{\mb q}^{(\sigma_2)i*} \nonumber \\
    &=(e_{\mb k_1}^{(\lambda_1)*}\cdot e_{\mb K}^{(\sigma_1)})
    (e_{\mb K}^{(\sigma_1)*}\cdot e_{\mb k_3}^{(\lambda_2)})
    (e_{\mb k_3}^{(\lambda_2)*}\cdot e_{\mb q}^{(\sigma_2)})
    (e_{\mb q}^{(\sigma_2)*}\cdot e_{\mb k_1}^{(\lambda_1)}).
\end{align}
Their explicit contraction check is collected in Appendix~B. The dimensionful loop integral is defined by
\begin{align}
    \widetilde{\mathcal{I}}^{\lambda_1\lambda_2}_{\sigma_1\sigma_2}
    &\equiv \be^{\si_1L} \be^{\si_2L} \int\frac{\di\mb q^3}{(2\pi)^3}
    \left(\frac{q}{2}\right)^{2\ii f\tn}
    \left(\frac{K}{2}\right)^{2\ii f\tn}
    \overline{\mathcal{E}}^{\lambda_1\lambda_2}_{\sigma_1\sigma_2} \nonumber \\
    &=\frac{4^{-2\ii f\tn}k_s^{3+4\ii f\tn}}{(4\pi)^{3/2}}
    \mathcal{I}^{\lambda_1\lambda_2}_{\sigma_1\sigma_2}(\theta_1,\theta_3,\phi_3),
    \qquad f=\pm.
\end{align}
Here the index $f$ labels the two nonlocal branches and is distinct from the SK index $d$ of the box
diagram. The loop integral can then be written in the same form as the triangle result:
\begin{align}
    \mathcal{LOOP}
    =\sum_{\lambda_1,\lambda_2=L,\pm}
    &\left[D^{(\lambda_1)}(k_1,\tau_1,\tau_2)\right]_{ab}
    \left[D^{(\lambda_2)}(k_3,\tau_3,\tau_4)\right]_{cd} \nonumber \\
    &\times\sum_{f=\pm}(\tau_1\tau_2\tau_3\tau_4)^{1/2+\ii f\tn}
    \mathcal{B}^{\lambda_1\lambda_2}_f(k_1,k_3),
\end{align}
where
\begin{align}
    \mathcal{B}^{\lambda_1\lambda_2}_f(k_1,k_3)
    &= \frac{\Gamma^4(-\ii f\tn)}{(4\pi)^{7/2}}
    4^{-2\ii f\tn}k_s^{3+4\ii f\tn}
    \sum_{\sigma_1,\sigma_2=L,\pm}
    \mathcal{I}^{\lambda_1\lambda_2}_{\sigma_1\sigma_2}(\theta_1,\theta_3,\phi_3) \nonumber \\
    &= \frac{\Gamma^4(-\ii f\tn)}{(4\pi)^{7/2}}
    4^{-2\ii f\tn}k_s^{3+4\ii f\tn}
    \mathcal{I}^{\lambda_1\lambda_2}(\theta_1,\theta_3,\phi_3)
\end{align}

Then we can deal with the left and right subgraph, which is identical to the left subgraph contribution of triangle diagrams:
\begin{align}
    \mathcal{T}^{(L),\lam_1}_f 
    &= \left( ig^2h_0\right)^2 
    \frac{k_{12}^{-2id\tn}}{4k_1^3 k_2^3} 
    \cdot \mathcal{P}_f^{(2),\lam_1}(\tn) \\
    \mathcal{T}^{(R),\lam_2}_f 
    &= \left( ig^2h_0\right)^2 
    \frac{k_{34}^{-2id\tn}}{4k_3^3 k_4^3} 
    \cdot \mathcal{P}_f^{(2),\lam_2}(\tn)
\end{align}
So the box diagram:
\begin{align}
    \la \delta h^4 \ra^{\text{Box}} 
    &= \sum_{f=\pm} \sum_{\lam_1,\lam_2=L,\pm} 
    \mc T^{(L),\lambda_1}_f  
    \cdot \mc B^{\lam_1\lam_2}_f 
    \cdot \mc T^{(R),\lam_2}_f \\
    &= (ig^2h_0)^4 
    \frac{H^2 k_s^3}{16 k_1^3 k_2^3 k_3^3 k_4^3} 
    \left( \frac{k_s^2}{4k_{12}k_{34}}\right)^{2\ii\tn} 
    \times C^{\text{Box}}(\tn,\theta_1) + \text{c.c.}
\end{align}
where the magnitude is:
\begin{align}
    C^{\text{Box}}(\tn,\theta_1,\theta_3,\phi_3) = 
    \sum_{\lam_z,\lam_2=L,\pm}
    [\mathcal{P}^{(2),\lam_1}_+(\tn) \mc P^{(2),\lambda_2}_+(\tn) ]
    \cdot \frac{\Gamma^4(-\ii \tn)}{(4\pi)^{7/2}} \cdot 
    \mathcal{I}^{\lambda_1\lam_2}(\theta_1,\theta_3,\phi_3)|_{f=+}
\end{align}


\appendix
\section{Check of the Box Polarization Contractions}

For the Box diagram, introduce the shorthand
\begin{align}
    \hat q&=(u\cos\phi,u\sin\phi,z),\qquad
    \hat K=\frac{1}{K}(qu\cos\phi,qu\sin\phi,qz+k_s), \\
    \hat k_1&=(s_1,0,c_1),\qquad
    \hat k_3=(s_3\cos\phi_3,s_3\sin\phi_3,c_3),
\end{align}
where $s_3=\sin\theta_3$ and $c_3=\cos\theta_3$. For a unit vector with polar angles
$(\theta_p,\phi_p)$, we use
\begin{align}
    e_{\mb p}^{(L)}&=(\sin\theta_p\cos\phi_p,\sin\theta_p\sin\phi_p,\cos\theta_p), \\
    e_{\mb p}^{(\sigma)}&=\frac{1}{\sqrt{2}}(\cos\theta_p\cos\phi_p-\sigma\ii\sin\phi_p,
    \cos\theta_p\sin\phi_p+\sigma\ii\cos\phi_p,-\sin\theta_p),
    \qquad \sigma=\pm1.
\end{align}
Thus $e_{\mb k_1}^{(\lambda_1)}$ and $e_{\mb k_3}^{(\lambda_2)}$ follow by taking
$(\theta_p,\phi_p)=(\theta_1,0)$ and $(\theta_3,\phi_3)$, respectively.

For the physical Box integral, the soft helicities must be kept explicit. Define
\begin{align}
    A^{\lambda_1}_{\sigma_1}&\equiv e_{\mb k_1}^{(\lambda_1)*}\cdot e_{\mb K}^{(\sigma_1)},
    &B^{\lambda_2}_{\sigma_1}&\equiv e_{\mb K}^{(\sigma_1)*}\cdot e_{\mb k_3}^{(\lambda_2)}, \\
    C^{\lambda_2}_{\sigma_2}&\equiv e_{\mb k_3}^{(\lambda_2)*}\cdot e_{\mb q}^{(\sigma_2)},
    &D^{\lambda_1}_{\sigma_2}&\equiv e_{\mb q}^{(\sigma_2)*}\cdot e_{\mb k_1}^{(\lambda_1)}.
\end{align}
The Box angular factor is therefore
\begin{align}
    \mathcal{E}^{\lambda_1\lambda_2}_{\sigma_1\sigma_2}
    =A^{\lambda_1}_{\sigma_1}B^{\lambda_2}_{\sigma_1}
    C^{\lambda_2}_{\sigma_2}D^{\lambda_1}_{\sigma_2},
\end{align}
which is the direct Box analogue of the three-factor Triangle contraction. One must first form this
product, then perform the $\phi$ average, and only afterwards apply the loop-seed map separately for
every $(\sigma_1,\sigma_2)$.

\iffalse
For reference, a signed polarization sum over the two soft propagators is represented by
\begin{align}
    P^{ij}(\hat p)&\equiv\sum_{\mu=L,\pm}\eta_\mu
    e_{\mb p}^{(\mu)i}e_{\mb p}^{(\mu)j*}
    =\delta^{ij}-2\hat p^i\hat p^j.
\end{align}
This identity must only be applied after the individual $(\mu,\nu)$ loop weights have been shown to be
identical. In the Box calculation we instead keep $(\mu,\nu)$ explicit, exactly as in the Triangle
calculation. The following $X,Y$ construction is therefore only an unweighted algebraic diagnostic and
is not used to define $\mathcal{I}^{\lambda_1\lambda_2}_{\mu\nu;f}$:
\begin{align}
    X_{\lambda_1\lambda_2}
    &\equiv e_{\mb k_1}^{(\lambda_1)*}\cdot P(\hat K)\cdot e_{\mb k_3}^{(\lambda_2)} 
    =e_{\mb k_1}^{(\lambda_1)*}\cdot e_{\mb k_3}^{(\lambda_2)}
    -2(e_{\mb k_1}^{(\lambda_1)*}\cdot\hat K)
    (\hat K\cdot e_{\mb k_3}^{(\lambda_2)}), \\
    Y_{\lambda_2\lambda_1}
    &\equiv e_{\mb k_3}^{(\lambda_2)*}\cdot P(\hat q)\cdot e_{\mb k_1}^{(\lambda_1)}
    =e_{\mb k_3}^{(\lambda_2)*}\cdot e_{\mb k_1}^{(\lambda_1)}
    -2(e_{\mb k_3}^{(\lambda_2)*}\cdot\hat q)
    (\hat q\cdot e_{\mb k_1}^{(\lambda_1)}).
\end{align}
The nine results are then
\begin{align}
    \left(\mathcal{E}^{\lambda_1\lambda_2}\right)_{\lambda_1,\lambda_2=L,+,-}
    =
    \begin{pmatrix}
    X_{LL}Y_{LL} & X_{L+}Y_{+L} & X_{L-}Y_{-L} \\
    X_{+L}Y_{L+} & X_{++}Y_{++} & X_{+-}Y_{-+} \\
    X_{-L}Y_{L-} & X_{-+}Y_{+-} & X_{--}Y_{--}
    \end{pmatrix}.
\end{align}
For equal internal weights, this expression is equivalent to the signed sum of the four-polarization
contractions, since
\begin{align}
    \mathcal{E}^{\lambda_1\lambda_2}
    =e_{\mb k_1}^{(\lambda_1)i}e_{\mb k_1}^{(\lambda_1)j*}
    P^{jk}(\hat K)e_{\mb k_3}^{(\lambda_2)k}e_{\mb k_3}^{(\lambda_2)l*}
    P^{li}(\hat q)
    =X_{\lambda_1\lambda_2}Y_{\lambda_2\lambda_1}.
\end{align}
It provides a short component-by-component diagnostic only. The physical Box coefficients are instead
obtained by first retaining the four scalar products in
$\mathcal{E}^{\lambda_1\lambda_2}_{\mu\nu}$ and then carrying out the $(\mu,\nu)$-dependent loop integrals.

For completeness, the azimuthal integral can be made fully explicit without listing several hundred
individual monomials. Let $a=e_{\mb k_1}^{(\lambda_1)}$, $b=e_{\mb k_3}^{(\lambda_2)}$, and denote their
transverse two-components by $a_\perp,b_\perp$. The six loop-seed combinations are
\begin{align}
    A_f={}&-2\mathcal{J}_0(0,2)+4\mathcal{J}_2(0,2)+\mathcal{J}_0(0,0)
    -2\mathcal{J}_2(0,0)-4\mathcal{J}_1(1,2)+8\mathcal{J}_3(1,2)
    -2\mathcal{J}_2(2,2)+4\mathcal{J}_4(2,2), \\
    B_f={}&\mathcal{J}_0(0,0)-2\mathcal{J}_2(0,0)-\mathcal{J}_0(2,2)
    +3\mathcal{J}_2(2,2)-2\mathcal{J}_4(2,2), \\
    C_f={}&-2\mathcal{J}_2(0,2)+\mathcal{J}_2(0,0)-4\mathcal{J}_3(1,2)
    -2\mathcal{J}_4(2,2), \\
    D_f={}&2\mathcal{J}_1(1,2)-2\mathcal{J}_3(1,2)
    +2\mathcal{J}_2(2,2)-2\mathcal{J}_4(2,2), \\
    F_f={}&\mathcal{J}_2(0,0)-\frac{1}{2}\mathcal{J}_0(2,2)
    +\frac{1}{2}\mathcal{J}_4(2,2), \\
    G_f={}&\frac{1}{2}\mathcal{J}_0(2,2)-\mathcal{J}_2(2,2)
    +\frac{1}{2}\mathcal{J}_4(2,2),
\end{align}
where every $\mathcal{J}_r$ is evaluated at $x=-\ii f\tn$. The $\mathcal{J}_4$ seed is reduced to
$\mathcal{J}_0$ by the closed forms given below.

The nine dimensionless coefficients are the following explicit direction functions:
\begin{align}
    \mathcal{I}^{\lambda_1\lambda_2}_f(k_1,k_3)
    ={}&A_f|a_z|^2|b_z|^2
    +B_fa_zb_z^*(a_\perp^*\cdot b_\perp)
    +C_fa_z^*b_z(a_\perp\cdot b_\perp^*) \nonumber \\
    &+D_f\left[a_zb_z(a_\perp^*\cdot b_\perp^*)
    +a_z^*b_z^*(a_\perp\cdot b_\perp)
    +|a_\perp|^2|b_z|^2+|a_z|^2|b_\perp|^2\right] \nonumber \\
    &+F_f|a_\perp^*\cdot b_\perp|^2
    +2G_f|a_\perp|^2|b_\perp|^2,
    \qquad (\lambda_1,\lambda_2)=L,+,-.
\end{align}
More explicitly, the nine entries are obtained by the substitution
\begin{align}
    \left(\mathcal{I}^{\lambda_1\lambda_2}_f\right)_{\lambda_1,\lambda_2=L,+,-}
    =\begin{pmatrix}
    \mathscr{F}_f(a_L,b_L)&\mathscr{F}_f(a_L,b_+)&\mathscr{F}_f(a_L,b_-)\\
    \mathscr{F}_f(a_+,b_L)&\mathscr{F}_f(a_+,b_+)&\mathscr{F}_f(a_+,b_-)\\
    \mathscr{F}_f(a_-,b_L)&\mathscr{F}_f(a_-,b_+)&\mathscr{F}_f(a_-,b_-)
    \end{pmatrix},
\end{align}
where $\mathscr{F}_f(a,b)$ denotes the preceding displayed direction function and
\begin{align}
    a_L&=(s_1,0;c_1),& a_\pm&=\frac{1}{\sqrt{2}}(c_1,\pm\ii;-s_1), \\
    b_L&=(s_3\cos\phi_3,s_3\sin\phi_3;c_3),&
    b_\pm&=\frac{1}{\sqrt{2}}(c_3\cos\phi_3\mp\ii\sin\phi_3,
    c_3\sin\phi_3\pm\ii\cos\phi_3;-s_3).
\end{align}
Here the semicolon separates transverse and $z$ components. Thus every entry depends only on
$\theta_1$, $\theta_3$, and $\phi_3$, and not on $k_1$ or $k_3$.

\fi

\section{Closed Gamma-function forms for $\mathcal{J}_3$ and $\mathcal{J}_4$}

For a general loop seed, write $\mathcal{J}_r(m,n)=f_r(a,b)$ with
\begin{align}
    a=x-\frac{m}{2},\qquad b=x+\frac{n}{2},\qquad x=-\ii f\tn.
\end{align}
The identity $z=(K^2-q^2-k_s^2)/(2qk_s)$ gives the following finite Gamma-function
closed forms:
\begin{align}
    f_3(a,b)
    =\frac{1}{8}\Big[&f_0\left(a+\frac{3}{2},b-3\right)
    -3f_0\left(a+\frac{3}{2},b-2\right)
    -3f_0\left(a+\frac{1}{2},b-2\right) \nonumber \\
    &+3f_0\left(a+\frac{3}{2},b-1\right)
    +6f_0\left(a+\frac{1}{2},b-1\right)
    +3f_0\left(a-\frac{1}{2},b-1\right) \nonumber \\
    &-f_0\left(a+\frac{3}{2},b\right)
    -3f_0\left(a+\frac{1}{2},b\right)
    -3f_0\left(a-\frac{1}{2},b\right)
    -f_0\left(a-\frac{3}{2},b\right)\Big], \\
    f_4(a,b)
    =\frac{1}{16}\Big[&f_0(a+2,b-4)
    -4f_0(a+2,b-3)-4f_0(a+1,b-3) \nonumber \\
    &+6f_0(a+2,b-2)+12f_0(a+1,b-2)+6f_0(a,b-2) \nonumber \\
    &-4f_0(a+2,b-1)-12f_0(a+1,b-1)-12f_0(a,b-1)-4f_0(a-1,b-1) \nonumber \\
    &+f_0(a+2,b)+4f_0(a+1,b)+6f_0(a,b)+4f_0(a-1,b)+f_0(a-2,b)\Big].
\end{align}
Here $f_0$ is the Gamma ratio defined in the loop-seed section. Consequently, these equations are
closed expressions in Gamma functions, not new independent seeds. More generally,
\begin{align}
    f_r(a,b)=\frac{1}{2^r}
    \sum_{p+q+t=r}\frac{r!}{p!q!t!}(-1)^{q+t}
    f_0\left(a+\frac{r}{2}-q,b-p\right),
\end{align}
which provides a direct check of both formulas.

\bibliography{sample} 
\bibliographystyle{utphys}

\end{document}
